\documentclass[11pt,a4paper]{article}
\usepackage[a4paper,margin=25mm]{geometry}
\usepackage[T1]{fontenc}
\usepackage{xcolor,lmodern,microtype,amsmath,amssymb,amsthm,booktabs}
\usepackage[colorlinks=true,linkcolor=black,citecolor=magenta,urlcolor=black]{hyperref}
\newtheorem{thm}{Theorem}[section]
\newtheorem{prop}[thm]{Proposition}
\newtheorem{lem}[thm]{Lemma}
\newtheorem{defn}[thm]{Definition}
\newtheorem{conj}[thm]{Conjecture}
\newtheorem*{introconj}{Conjecture}
\newtheorem{introthm}{Theorem}
\renewcommand{\theintrothm}{\Alph{introthm}}
\setlength{\parskip}{1pt}
\linespread{0.99}
\setlength{\emergencystretch}{2em}
\setcounter{tocdepth}{1}
\raggedbottom
\title{Numerical evidence for the almost Hadamard conjecture}
\author{Chat, Grok and Teo}
\date{}
\hypersetup{pdftitle={Numerical evidence for the almost Hadamard conjecture},pdfauthor={Chat, Grok and Teo}}
\begin{document}
\maketitle
\begin{abstract}
We report on encouraging numerical evidence for the almost Hadamard conjecture, stating that any local maximizer of the $1$-norm on $U_N$ must have all entries of modulus $1/\sqrt{N}$, or equivalently that every complex almost Hadamard matrix $H=\sqrt{N}\,U$ must be a complex Hadamard matrix, with computations at the key values $N=4,5$ arising from our attempts to both prove and disprove the conjecture, and adding to the previous abstract evidence.
\end{abstract}
\begingroup
\color{magenta}\hypersetup{linkcolor=magenta}
\setlength{\parskip}{0pt}
\renewcommand{\addvspace}[1]{\vspace{3pt}}
\tableofcontents
\endgroup
\section*{Introduction}
\phantomsection\addcontentsline{toc}{section}{Introduction}
A complex Hadamard matrix is a square matrix whose entries have modulus
1 and whose rows are orthogonal:
\[
 |H_{ij}|=1,\qquad HH^*=NI
\]
Its rescaling $U=H/\sqrt N$ is a flat unitary, meaning that all its
entries have modulus $1/\sqrt{N}$. These matrices attain the upper bound
\[
 L(U)=\|U\|=\sum_{ij}|U_{ij}|\le N\sqrt N
\]
The following conjecture was made in \cite{BN}.

\begin{introconj}[Almost Hadamard, AHC]
Every local maximizer of the $1$-norm on $U_N$ attains the bound
$N\sqrt N$.
\end{introconj}

The starting point was the real theory \cite{BCS,BNZ}. A real almost
Hadamard matrix is a rescaled local maximizer on $O_N$; such matrices
exist in every order and need not be Hadamard. The passage to $U_N$
allows new directions of variation. These exclude several real
families coming from designs and circulants \cite{BN}, as recalled in
Section~\ref{background}.

\begin{samepage}
In order to put this conjecture in context and illustrate its difficulty,
we discuss 2 other conjectures concerning unitary groups.

\begin{introconj}[Marcus--de Oliveira]
For diagonal matrices $A=diag(a_1,\ldots,a_N)$ and
$B=diag(b_1,\ldots,b_N)$, the inclusion
\[
 \det(A+UBU^*)\in conv
 \left\{\prod_{i=1}^N(a_i+b_{\sigma(i)}):\sigma\in S_N\right\}
\]
holds for every $U\in U_N$.
\end{introconj}
\end{samepage}
This concerns sums of normal matrices with prescribed spectra
\cite{M,O}: unitary diagonalization gives the displayed form. The points
on the right pair the eigenvalues in every possible order. The
conjecture confines the determinant to their convex hull in the complex
plane. It is known through $N=3$ \cite{BMP}; $N=4$ is the first
unsettled order \cite{BP,MM}.

Higher-dimensional results treat self-adjoint pairs \cite{F}, essentially
self-adjoint pairs \cite{BKP}, scalar multiples of unitaries \cite{D}, and
certain normal dilations \cite{R}. A variational approach maximizes
the projection of the determinant onto a line and studies the resulting
stationarity equations on $U_N$ \cite{PB,BP}. Its successes and
obstructions \cite{MM} provide a useful comparison with the AHC.

The second conjecture concerns squared entry moduli. For a unitary $U$,
put
\[
 p(U)=(|U_{ij}|^2),\qquad J_N=(1/N)_{ij}
\]
The matrix $p(U)$ is bistochastic: its entries are $\ge0$ and its
row and column sums are 1. Bistochastic matrices arising this way
are called unistochastic.

\begin{introconj}[Unistochastic star-shapedness]
For every unistochastic matrix $M$, the segment
\[
 M_t=(1-t)M+tJ_N,\qquad 0\le t\le1
\]
consists of unistochastic matrices.
\end{introconj}
This holds at $N=3$ \cite{BEKTZ}. Rajchel, G\k{a}siorowski and
\.Zyczkowski proposed the order-4 conjecture in \cite{RGZ}; the
higher-dimensional version is discussed in \cite{RKKZ}. Partial
results include rays supplied by robust Hadamard matrices \cite{RGZ}
and the circulant subset at order 4 \cite{RKKZ}. In every order,
the larger set defined by the necessary row and column polygon
inequalities, called bracelet conditions, is star-shaped \cite{RKKZ}.
This does not settle the unistochastic case.

The connection to the AHC is the strict increase
\[
 h(M_t)>h(M)\quad(M\ne J_N,\ t>0),\qquad
 h(M)=\sum_{ij}\sqrt{M_{ij}},\qquad L=h\circ p
\]
To contradict a local maximum at a particular $U$, however, one needs
unitaries realizing $M_t$ that approach that same $U$, up to phases.
Section~\ref{orderthree} supplies this lifting argument at order 3.
Reconstruction from moduli \cite{AMM}, the defect \cite{TZ}, and the
critical loci of the modulus map \cite{BebN} describe the higher-order
difficulty. Random unistochastic ensembles \cite{ZKSS}, volume
calculations \cite{DZ}, manifold optimization \cite{AMS}, and calibrated
searches for mutually unbiased bases \cite{BW} give further context.

Our principal structural results concern $U_4$.

\begin{introthm}
If one $2\times2$ submatrix of a local maximum on $U_4$ can be made
real by row and column phases, the whole unitary is flat.
\end{introthm}

The proof uses a block variation with positive second derivative.
When no $2\times2$ block can be made real by phases, the critical-point
equations instead impose restrictions on the entry moduli.

\begin{introthm}
A local maximum on $U_4$ is flat if every row and every column contains
the same multiset of entry moduli.
\end{introthm}

In the remaining case, the reciprocal moduli satisfy a rank-2
condition and explicit polynomial identities. Combining these
restrictions with a count of the possible modulus patterns gives
another criterion.

\begin{introthm}
A nonflat local maximum on $U_4$ would have at least $5$ distinct
entry moduli.
\end{introthm}

The last proof uses an exact enumeration of $27$ patterns. All 3
proofs are in Section~\ref{orderfour}, with the remaining curvature
equations and numerical evidence in Section~\ref{fourcurvature}.

At order 5, our explicit saddle has
\[
 \text{Hessian inertia}=(15,9,1),\qquad
 \text{all row and column pair curvatures}<0
\]
The inertia lists negative, zero and positive eigenvalues. A
3-column variation gives an increase, and both unstable gradient
branches converge to Fourier matrices (Section~\ref{fivegeometry}).
Thus pair tests alone can miss a saddle.

The original campaigns and the subsequent targeted search have the
following sizes, each from a 7-hour run:
\begin{center}
\begin{tabular}{clrr}
\toprule
Order & Campaign & Ascent starts & Stationary-search starts\\
\midrule
$4$ & Original & $83\,076\,896$ & $52\,237\,424$\\
$5$ & Original & $181\,164\,464$ & $22\,933\,376$\\
$5$ & Targeted & $8\,737\,840$ & $26\,213\,552$\\
\bottomrule
\end{tabular}
\end{center}
At order 4 every ascent return is within $6.014\cdot10^{-6}$ of
flat in entry moduli. Every accepted nonflat stationary return has
positive curvature; the near-zero and unresolved returns lie below a necessary
entry bound. At order 5 every ascent return is near flat, while
the original stationary campaign left $404\,581$ unresolved returns. Among its saved
returns, $452$ nonflat stationary matrices form $40$ numerical groups
and have interval exclusions. Each group representative has a nearby
exact critical point with certified Hessian inertia. Their norm
enclosures give at least $35$ distinct nonflat critical values.
All $288$ saved stalls are explained.
One explanation is a certified positive local minimum of the squared
stationarity residual: minimizing that residual need not find a
critical point of $L$.

The targeted campaign leaves no unresolved or candidate flags under
the endpoint criteria. A separate replay of the first $16\,189\,008$
original starts, about $70.6\%$ of that sequence, also leaves no such
flags after follow-up. Its remaining $6\,744\,368$ starts were not
revisited, and the replay is not added to the table. Small-entry
exclusions concern the computed endpoints, not their possible limits;
matching the known invariant signatures does not prove equivalence.

These are sampled searches, not coverings of $U_4$ or $U_5$.
Sections~\ref{framework} and~\ref{fivecampaign} separate the variational
criteria from the numerical classifications and certificates.
From \texttt{code011}, with NumPy, SciPy, mpmath and SymPy,
\texttt{python3 reproduce.py}
rechecks the identities, certificates and retained outputs in a
temporary copy. Its report \texttt{data/reproduction.json} records
versions, hashes and outputs for all $27$ checks.
\texttt{README.md} indexes the code, and every program is cited below.

\section{Hadamard matrices and the conjecture}\label{background}
\subsection{Hadamard matrices and norm maximization}
The symbols $*$ and $t$ denote conjugate transpose and transpose.

\begin{samepage}
\begin{defn}
A complex Hadamard matrix of order $N$ is a matrix $H$ satisfying
\[
 |H_{ij}|=1,\qquad HH^*=NI
\]
It is real Hadamard when its entries belong to $\{1,-1\}$.
A unitary $U$ is flat when $|U_{ij}|=1/\sqrt{N}$ for every entry.
\end{defn}
\end{samepage}

The Fourier matrices provide complex examples in every order:
\[
 F_N=(w^{ij})_{0\le i,j<N},\qquad w=e^{2\pi i/N}
\]
For instance, with $w=e^{2\pi i/3}$ in the second matrix,
\[
 F_2=\begin{pmatrix}1&1\\1&-1\end{pmatrix},\qquad
 F_3=\begin{pmatrix}1&1&1\\1&w&w^2\\1&w^2&w\end{pmatrix}
\]
Their row orthogonality follows from the geometric-series identity.
For real Hadamard matrices, $N\ge3$ requires $4$ to divide $N$.
See \cite{TZguide} for background.

\begin{prop}\label{globalbound}
The entrywise $1$-norm satisfies
\[
 L(U)=\|U\|=\sum_{ij}|U_{ij}|\le N\sqrt N\qquad(U\in U_N)
\]
Equality holds exactly when $U$ is flat, or equivalently when
$H=\sqrt N\,U$ is complex Hadamard.
\end{prop}
\begin{proof}
Cauchy--Schwarz and unitarity give
\[
 \sum_{ij}|U_{ij}|\le N\Big(\sum_{ij}|U_{ij}|^2\Big)^{1/2}=N\sqrt N
\]
Equality means that all entry moduli are equal. Their squares sum to
$N$, so their common value is $1/\sqrt{N}$.
\end{proof}

\subsection{Almost Hadamard matrices}
The real theory replaces the equality case above by a local maximum
\cite{BCS,BNZ}. The complex extension is introduced in \cite{BN}.

\begin{defn}
Write $H=\sqrt N\,U$. The matrix $H$ is real almost Hadamard if $U$
locally maximizes $L$ on $O_N$, and complex almost Hadamard if $U$
locally maximizes $L$ on $U_N$.
\end{defn}

Compactness gives real almost Hadamard matrices in every order.
``Almost'' refers to local maximality, with no prescribed distance
from a Hadamard matrix. The simplest nonflat examples are reflections.

\begin{prop}\label{realreflection}
For $N\ge3$, the matrix
\[
 K_N=\frac1{\sqrt N}
 \begin{pmatrix}
 2-N&2&\cdots&2\\
 2&2-N&\cdots&2\\
 \vdots&\vdots&\ddots&\vdots\\
 2&2&\cdots&2-N
 \end{pmatrix}
\]
is real almost Hadamard. It is Hadamard exactly when $N=4$.
\end{prop}
\begin{proof}
Let $J_N=(1/N)_{ij}$ and $W_N=K_N/\sqrt N=2J_N-I$.
Since $J_N^2=J_N$, the matrix $W_N$ is orthogonal. For its sign matrix
$S$, direct multiplication gives
\[
 S^tW_N=2I+(N-4)J_N>0
\]
Near $W_N$ the entry signs are fixed and $L(V)=Tr(S^tV)$.
For every nonzero real skew-symmetric $A$, the curve $W_Ne^{tA}$ has
\[
 L'(0)=0,\qquad L''(0)=Tr((S^tW_N)A^2)<0
\]
Thus $W_N$ is a strict local maximum on $O_N$. The diagonal and
off-diagonal moduli agree exactly when $N-2=2$.
\end{proof}

The same matrices fail to be complex almost Hadamard when $N\ne4$
\cite[Theorem 4.7]{BN}. This is the basic real--complex distinction:
the additional unitary directions can destroy a real local maximum.

\begin{conj}[Almost Hadamard conjecture]\label{ahc}
Every local maximizer of the $1$-norm on $U_N$ is flat. Equivalently, every
complex almost Hadamard matrix is complex Hadamard.
\end{conj}

Local maximality here is unrestricted and need not be strict. Row and
column phases preserve $L$, and Hadamard matrices can occur in families.
The real theory also includes circulant and 2-entry constructions
\cite{BNZ}, other entrywise exponents \cite{BNp}, and polar factors of
complementary Hadamard submatrices under suitable size assumptions
\cite{BNS}.

\subsection{Designs and circulants: the earlier evidence}
The 2-entry constructions come from square incidence patterns of
symmetric block designs. Their rows have constant weights and constant
pairwise overlaps.

\begin{defn}
An $(a,b,c)$ pattern has $2$ symbols $x,y$ such that every pair of
distinct rows has the following counts:
\[
 \begin{array}{c|cccc}
 \text{pair}&(x,x)&(x,y)&(y,x)&(y,y)\\
 \hline
 \text{count}&a&b&b&c
 \end{array}
 \qquad N=a+2b+c
\]
We take $a,c\ge0$ and $b>0$.
\end{defn}

For a full-support unitary with this pattern, normalize $y>0$ and write
$x=-\varepsilon t y$, with $t>0$ and $|\varepsilon|=1$.
Orthogonality and normalization become
\[
 at^2-2b\,Re(\varepsilon)t+c=0,\qquad
 y=\big((a+b)t^2+b+c\big)^{-1/2}
\]
Criticality further restricts the phase $\varepsilon$.

\begin{prop}
A full-support critical unitary with an $(a,b,c)$ pattern is, after
multiplication by a scalar of modulus $1$, real or flat.
\end{prop}
\begin{proof}
The criticality calculation of \cite[Theorem 5.2]{BN} gives
\[
 (\varepsilon-\bar\varepsilon)(t-1)=0
\]
If $t=1$ the 2 entry moduli agree. Otherwise $\varepsilon$ is real,
and orthogonality excludes $\varepsilon=-1$, leaving $\varepsilon=1$.
\end{proof}

The real branches still require a second-variation test.
Theorem 6.5 of \cite{BN} excludes $K_N$ for $N\ne4$, the Grassmannian
matrices $I_N^{q,d}$, and the Paley biplane matrix $P_{11}$.
The projective-plane examples, including the Fano plane at $q=2$,
belong to the Grassmannian family at $d=1$ \cite[Proposition 5.5]{BN}.
These exclusions do not cover every construction from a design.

A circulant matrix has rows obtained by successive cyclic shifts. Its
unitarity can be expressed in Fourier coordinates as
\[
 U=\frac1N F_N\,diag(q_0,\ldots,q_{N-1})F_N^*,\qquad |q_k|=1
\]
For real symmetric circulants the parameters satisfy
\[
 q_k\in\{1,-1\},\qquad q_k=q_{-k}
\]
The method of \cite[Section 7]{BN} averages second derivatives over
suitable circulant self-adjoint generators. A positive average supplies
an increasing direction. It excludes the explicit nonflat circulant almost
Hadamard examples of \cite{BNZ}; further computer checks use the
examples of \cite{BNp}. A general averaged inequality is proposed in
\cite[Conjecture 7.8]{BN}, for real symmetric circulants and, in an
analogous form, for self-adjoint unitary circulants. It is not proved there.

This earlier evidence gives families and variation methods. The
small-order study below tests the full tangent space; the order-5
model shows why restricted tests can miss an increasing direction.

\section{Variations and small orders}\label{framework}
We now pass from norm maximization to first and second derivatives.
The only nonsmooth points are matrices with zero entries.

\subsection{Criticality and the Hessian}
\begin{lem}\label{fullsupport}
A local maximum of $L$ on $U_N$ has no zero entry.
\end{lem}
\begin{proof}
Choose a column rotation that moves a nonzero entry into a zero
position to first order. In the sum of the increments at $t$ and $-t$,
nonzero entries contribute $O(t^2)$, the chosen zero contributes a
positive multiple of $|t|$, and the other zeros contribute amounts $\ge0$.
This contradicts local maximality. See also \cite[Theorem 1.6]{BN}.
\end{proof}

We can therefore work where the norm is smooth. Its derivatives are
expressed in terms of the entry phases.

\begin{defn}
A matrix has full support when none of its entries is zero. For such
a unitary $V$, its phase matrix and associated matrix are
\[
 S_{ij}=\frac{V_{ij}}{|V_{ij}|},\qquad X=S^*V
\]
The matrix $V$ is critical for $L$ when the first derivative vanishes
along every unitary curve through $V$.
\end{defn}

We use the real Frobenius inner product and norm
\[
 \langle A,B\rangle=Re\,Tr(A^*B),\qquad
 \|A\|_F^2=\sum_{ij}|A_{ij}|^2
\]
Every tangent vector at $V$ is realized by a curve
\[
 V(t)=Ve^{itB},\qquad V'(0)=iVB,\qquad B=B^*
\]
Thus self-adjoint matrices $B$ parametrize all unitary directions.

\begin{lem}\label{variations}
Along $V e^{itB}$ the first derivative of $L$ at zero is
$-Im\,Tr(XB)$. The second derivative is
\[
 Q_V(B)=\sum_{ij}\frac{Re((VB)_{ij}\bar S_{ij})^2}
 {|V_{ij}|}-Re\,Tr(XB^2)
\]
Thus $V$ is critical precisely when $X$ is self-adjoint.
\end{lem}
\begin{proof}
For one entry $z(t)$ put $s=z(0)/|z(0)|$. Differentiating $|z|^2=z\bar z$
gives $(|z|)'=Re(z'\bar s)$ at zero, then
$(|z|)''=Re(z''\bar s)+Im(z'\bar s)^2/|z(0)|$.
Substitute $z'=i(VB)_{ij}$ and $z''=-(VB^2)_{ij}$, and sum. Testing
the first derivative against all self-adjoint $B$ gives the last assertion.
\end{proof}

At a critical point, $Q_V$ is the Hessian. Local maximality requires
\[
 Q_V(B)\le0\qquad\text{for every self-adjoint }B
\]
One positive value excludes a local maximum. Nonpositivity alone need
not establish one, since higher-order terms can matter.

\begin{prop}\label{phasequotient}
At full support the row and column phase orbits have dimension $2N-1$.
At a critical point their tangent directions lie in the kernel of the
Hessian, so the Hessian descends to a quotient of dimension $(N-1)^2$.
\end{prop}
\begin{proof}
The action is $V\to D_1VD_2$, with $D_1,D_2$ diagonal unitaries.
At full support it fixes $V$ exactly when $D_1=zI$ and $D_2=z^{-1}I$.
The effective action is free and its orbits have dimension $2N-1$.
Phase invariance makes the derivative of $L$ along every phase vector
field identically zero. Differentiating at a critical point shows that
all Hessian terms involving a phase direction vanish. Removing these
directions from the $N^2$-dimensional tangent space gives the result.
\end{proof}

The horizontal Hessian is the restriction to the orthogonal complement
of the phase directions. The dimensions used in our computations are
\begin{center}
\begin{tabular}{lccc}
\toprule
Space & General order & $N=4$ & $N=5$\\
\midrule
Full tangent space & $N^2$ & $16$ & $25$\\
Phase directions & $2N-1$ & $7$ & $9$\\
Horizontal space & $(N-1)^2$ & $9$ & $16$\\
\bottomrule
\end{tabular}
\end{center}
A positive geodesic second derivative also excludes a local maximum
without requiring an exact stationarity certificate.

\subsection{Pair tests and a polar condition}
For 2 columns write their entries and relative phases as
\[
 p_i u_i,\quad q_i v_i,\qquad
 p_i,q_i>0,\quad |u_i|=|v_i|=1,\qquad h_i=\bar u_i v_i
\]
After removing column phases there are 2 real mixing directions.
Put
\[
 d_i=\frac{q_i^2}{p_i}+\frac{p_i^2}{q_i}
\]
These quantities determine the largest pair curvature.

\begin{lem}\label{twocol}
For $2$ columns as above, the quantity
\[
 M_{jk}=\frac14\left(\sum_i d_i+
       \left|\sum_i d_i h_i^2\right|\right)-\frac12\sum_i(p_i+q_i)
\]
is the largest value of $Q_V$ on Frobenius-unit off-diagonal self-adjoint
generators supported on these columns.
\end{lem}
\begin{proof}
Set $B_{jk}=e^{i\theta}/\sqrt2$, $B_{kj}=\bar B_{jk}$. The first term
of $Q_V(B)$ in Lemma~\ref{variations} becomes $\frac12\sum_i d_i\,Re(h_i e^{-i\theta})^2$.
The second uses $B^2=(E_{jj}+E_{kk})/2$. Now
$2Re(z)^2=1+Re(z^2)$ on the unit circle.
Maximizing the resulting expression over $\theta$ gives the stated formula for $M_{jk}$.
\end{proof}

The same test applies to rows. Pair tests are inexpensive but do not
exhaust the Hessian, as the order-5 example will show. Rank-1
generators give a further necessary condition.

\begin{prop}
At a local maximum on $U_N$ we have $S^*U\ge I/N$. In particular $S$
is invertible and its unitary polar factor is $U$.
\end{prop}
\begin{proof}
Write $r_{ij}=|U_{ij}|$ and $X=S^*U$. Fix a unit vector $v$ and insert $B=vv^*$ into $Q_U(B)\le0$.
For $a_{ij}=Re((Uv)_i\bar v_j\bar S_{ij})$ we have
$\sum_j r_{ij}a_{ij}=|(Uv)_i|^2$. By Cauchy--Schwarz and
$\sum_j r_{ij}^3\le\sum_j r_{ij}^2=1$ we obtain
\[
 v^*Xv\ge\sum_{ij}\frac{a_{ij}^2}{r_{ij}}
 \ge\sum_i\frac{|(Uv)_i|^4}{\sum_j r_{ij}^3}
 \ge\sum_i |(Uv)_i|^4\ge\frac1N
\]
Criticality gives $X=X^*=U^*S$, hence $S=UX$ is the polar decomposition.
\end{proof}

\subsection{The modulus map and its defect}
\begin{defn}
A matrix $M$ is bistochastic when
\[
 M_{ij}\ge0,\qquad \sum_jM_{ij}=1,\qquad\sum_iM_{ij}=1
\]
It is unistochastic when $M=p(U)$ for a unitary $U$, where
\[
 p(U)=(|U_{ij}|^2)
\]
The fiber over $M$ consists of all unitaries with these squared moduli.
\end{defn}

The modulus map forgets phases. A tangent variation of a bistochastic
matrix has zero row and column sums, giving $(N-1)^2$ independent real
coordinates. The differential $dp_U$ records which variations can be
realized to first order by a curve of unitaries.

\begin{defn}
The defect of a full-support unitary is
\[
 d(U)=(N-1)^2-rank(dp_U)
\]
It measures the missing rank of the modulus map.
\end{defn}

This is the defect of \cite{TZ}. Nonflat critical points must lie in
the singular part of the modulus map:

\begin{prop}\label{defect}
At a full-support nonflat critical point of $L$ on $U_N$, the modulus
map $p(U)=(|U_{ij}|^2)$ has rank strictly less than $(N-1)^2$.
\end{prop}
\begin{proof}
The affine space of bistochastic matrices has dimension $(N-1)^2$.
Write $M=p(U)$ and $h(M)=\sum_{ij}\sqrt{M_{ij}}$. Concavity gives
\[
 dh_M(J_N-M)\ge h(J_N)-h(M)>0
\]
At a critical point of $L=h\circ p$, the differential $dh_M$ vanishes
on the image of $dp_U$. That image is therefore a proper subspace of
the space of matrices with zero row and column sums.
\end{proof}

A direction in the kernel of $dp_U$ leaves the squared moduli fixed to
first order. Along a unitary curve with this initial direction,
$p'(0)=0$ gives
\[
 L''(0)=\frac12\sum_{ij}\frac{p''_{ij}(0)}{|U_{ij}|}
\]
At a critical point, this second variation depends only on the initial tangent vector.
Finding a positive value in this kernel, after removing phases, is one
route to the AHC. Section~\ref{fourphases} implements it at order 4.
Averaging over the unit sphere of the horizontal kernel is another
possible test, following \cite{BN}; a generally positive average is not
established here.

\subsection{Orders 1 to 3 and the lifting argument}\label{orderthree}
\begin{defn}
A full-support unitary is dephased when its first row and column are
positive real. A path of unitaries $U(t)$ satisfying
\[
 p(U(t))=M(t)
\]
is called a lift of the modulus path $M(t)$.
\end{defn}

Row and column phases always give a dephased representative. The local
lifting problem asks for representatives approaching the specified
unitary, rather than another point of its fiber.

\begin{thm}\label{threeflat}
Every local maximum of $L$ on $U_N$ is flat for $N\le3$.
\end{thm}
\begin{proof}
At order 1 the assertion is automatic. At order 2, full support
and dephasing reduce the unitary to
\[
 \begin{pmatrix}\cos t&\sin t\\\sin t&-\cos t\end{pmatrix},\qquad
 0<t<\pi/2
\]
Its norm $2(\cos t+\sin t)$ has its only local maximum at $t=\pi/4$.

At order 3 a local maximum $U$ has full support. Multiply its rows and columns by
phases so that the first row and column are positive. At order 3,
the unistochastic set is star-shaped about $J_3$, and a positive
unistochastic matrix determines its dephased unitary up to complex
conjugation \cite{BEKTZ,AMM}.

Suppose $M=p(U)\ne J_3$. Choose a decreasing sequence $t_n\to0$ and dephased unitary
lifts $V_n$ of $(1-t_n)M+t_nJ_3$. Compactness of $U_3$ gives a convergent
subsequence. Its limit lies over $M$ and is still dephased, since every
entry of $M$ is positive. The limit is therefore $U$ or $\bar U$.
Conjugating the subsequence in the second case gives unitaries tending
to $U$. Concavity and $h(J_3)>h(M)$ yield
\[
 L(V_n)=h((1-t_n)M+t_nJ_3)
 \ge (1-t_n)h(M)+t_n h(J_3)>L(U)
\]
This contradicts local maximality.
\end{proof}

The argument needs both a segment in the image of $p$ and control of
its lifts. Higher-order fibers can be more complicated: star-shapedness
alone would leave the lifting question at a prescribed unitary. The
direct variations below avoid this assumption.

\section{Order 4: structural restrictions}\label{orderfour}
We first prove Theorem A, then obtain equations for the remaining phase
patterns and deduce Theorems B and C. Section~\ref{fourcurvature} then treats curvature and the numerical evidence. Throughout, $U$ has full support and
$S_{ij}=U_{ij}/|U_{ij}|$.

\subsection{A real submatrix forces flatness}\label{fourphases}
For a column pair $p_iu_i,q_iv_i$, call the real lines through the
phase quotients $h_i=\bar u_iv_i$ its phase lines; $h$ and $-h$ count
as 1 line. For a full-support $2\times2$ block put
\[
 j(E)=Im(\bar E_{11}E_{12}E_{21}\bar E_{22})
\]
Thus $j(E)=0$ exactly when phases make $E$ real.
\begin{thm}\label{threecase}
A local maximum on $U_4$ is flat if one column pair uses at most $3$
phase lines. Equivalently, every nonflat local maximum has $j(E)\ne0$
for all $36$ $2\times2$ submatrices.
\end{thm}
We separate the possibilities of 1, 2 and 3 phase lines.
The basic variation is an imaginary rotation of 2 real columns.

\begin{lem}\label{realup}
A real unitary locally maximizing $L$ on the complex unitary group is flat.
\end{lem}
\begin{proof}
Mix 2 columns by diagonal entries $\cos t$ and off-diagonal entries
$i\sin t$. Their squared moduli $a^2,b^2$ in each row become
$(1-s)a^2+sb^2,sa^2+(1-s)b^2$, where $s=\sin^2t$.
Strict concavity of the square root increases their sum unless
$a^2=b^2$. At a maximum every row therefore has constant moduli.
\end{proof}
The same rotation handles a pair whose relative phases lie on a
single real line. The remaining columns determine the equality case.

\begin{lem}\label{onepair}
A nonflat full-support $U\in U_4$ with a collinear phase-quotient pair
has an increasing arbitrarily small $2$-row or $2$-column rotation.
\end{lem}
\begin{proof}
Make the 2 columns real. Lemma~\ref{realup}'s rotation works unless
they are $(a_i),(\epsilon_i a_i)$ with $a_i>0$, $\epsilon_i=\pm1$.
Each sign has squared mass $1/2$ and must occur twice: a singleton would
force zeros in the remaining columns. Reorder to obtain
$(a,b,c,d)^t,(a,b,-c,-d)^t$, with $a^2+b^2=c^2+d^2=1/2$.
The remaining columns are $(e,f)W$, $W\in U_2$, where
\[
 e=\sqrt2(b,-a,0,0)^t,\qquad f=\sqrt2(0,0,d,-c)^t
\]
The first 2 rows are collinear in phase, as are the last 2.
The same rotation increases $L$ unless $a=b$ and $c=d$, making all
4 numbers $1/2$. Orthogonality of $W$ makes the last 2 column
quotients collinear. Mixing them increases $L$ unless their moduli
also agree, which forces flatness by the row norms. Transposition
gives the row version.
\end{proof}
At a critical point, 2 phase lines reduce to the preceding 1-line
case through orthogonality and the first-order equations.

\begin{lem}\label{twosingle}
At a full-support critical point on $U_4$, a row pair with exactly $2$
phase lines forces a collinear column pair, and conversely by transposition.
\end{lem}
\begin{proof}
The equations $\sum p_iq_ih_i=\sum(p_i-q_i)h_i=0$ split on the 2
lines, each containing 2 opposite phases. On each, products and
differences agree, determining the positive moduli. Rephase the rows to
\[
 (a,a,b,b),\qquad(c,-c,dz,-dz),\qquad z\notin\mathbb R
\]
If columns 1 and 2 differ in modulus on the bottom rows, their 2
remaining differences have opposite signs. Stationarity makes their
bottom phase quotients equal; orthogonality then makes these real,
giving collinearity. The same applies to columns 3 and 4.
Otherwise each bottom row has equal moduli within both pairs and lies
in the span of $(b,b,-a,-a)$ and $(d,-d,-cz,cz)$.
For its coefficients $\lambda,\mu$, the modulus equalities imply
$Re(\bar\lambda\mu)=Re(z\bar\lambda\mu)=0$, so
$\lambda\mu=0$. Columns 1 and 2 are again collinear.
\end{proof}
3 phase lines require a block variation. We first identify 8
independent modulus coordinates for the differential of the modulus map.

\begin{lem}\label{eightblocks}
Write $U=\left(\begin{smallmatrix}A&B\\C&D\end{smallmatrix}\right)$
in $2\times2$ blocks. At full support, $j(B)j(C)\ne0$ makes $dp_U$
surjective onto the $4$ coordinates of $|A|^2$, the first column of
$|B|^2$ and the first row of $|C|^2$.
\end{lem}
\begin{proof}
The block $C$ is invertible, since rank 1 would give $j(C)=0$.
The left skew-adjoint generator
$\left(\begin{smallmatrix}0&Z\\-Z^*&0\end{smallmatrix}\right)$
gives $\dot A=ZC$, prescribing the first 4 coordinates.
Keeping $A$ fixed, a last-column mixing changes the 2 entries of
$|B_{i1}|^2$ by $2Re(w\bar B_{i1}B_{i2})$ up to sign. These forms
are independent precisely when $j(B)\ne0$, and $C$ stays fixed.
Mixing the last rows likewise prescribes the 2 $C$ coordinates.
\end{proof}
At a nonflat critical point, Proposition~\ref{defect} bounds the rank
by 8. Under this lemma it is exactly 8, and these coordinates
determine every vector in its image. Put $M=p(U)$ and $h(M)=\sum\sqrt{M_{ij}}$.
\begin{lem}\label{blockcurve}
Let $U$ be a full-support nonflat critical point with $j(A)=0$ and
$j(B)j(C)\ne0$. Simultaneously exchange rows $1,2$ and
$3,4$ and likewise columns, obtaining $\tilde M$.
There is a unitary curve with $p'(0)=0$, $j(D)=0$, and
\[
 L''(0)=\sum_{\{x,y\}}\frac{(x-y)^2(x+y)}{xy}
\]
where the $8$ pairs $\{x,y\}$ are the exchanged entry moduli.
\end{lem}
\begin{proof}
Make $A$ real. For $E^\#=\left(\begin{smallmatrix}
\bar E_{22}&-\bar E_{21}\\-\bar E_{12}&\bar E_{11}
\end{smallmatrix}\right)$, the operation respects products and adjoints,
and $V^\#=\overline{\det V}V$ for $V\in U_2$.
Polar decomposition and unitarity give fixed $V,W\in U_2$ with
\[
 B=(I-AA^*)^{1/2}V,\quad C=W(I-A^*A)^{1/2},\quad D=-WA^*V
\]
Conversely these formulas complete any strict contraction $A$ to a
unitary. Use $A(t)=\cos t\,A+i\sin t\,A^\#$.
The products $A(t)^*A(t),A(t)A(t)^*$ are even, since the cross products
for real $A$ equal $(\det A)I$. Thus $B'(0)=C'(0)=0$, while
\[
 D(t)=\cos t\,D-i\sin t\,\eta D^\#,\qquad \eta=\det W\det V
\]
The 8 coordinates above have zero derivative; hence all of $p'(0)$
vanishes. In particular $\eta\overline{D_{11}D_{22}}$ and
$\eta\overline{D_{12}D_{21}}$ are real, so $j(D)=0$.
In the diagonal blocks, squared moduli follow exactly
$\cos^2t\,M+\sin^2t\,\tilde M$.

The difference $p''(0)-2(\tilde M-M)$ vanishes there and has zero
row and column sums. Each other block is a multiple of
$\left(\begin{smallmatrix}1&-1\\-1&1\end{smallmatrix}\right)$.
Both variations lie in $im(dp_U)$: since $j(D)=0$, some
last-column mixing fixes $|D|^2$ to first order, while $j(B)\ne0$
makes its change in $|B|^2$ nonzero. Transpose for $C$.
Criticality kills these differences under $dh_M$; consequently
$L''(0)=2dh_M(\tilde M-M)$, giving the formula by pairing terms.
\end{proof}
The curvature formula is positive unless all exchanged moduli agree.
We now rule out this equality case at a nonflat critical point.

\begin{lem}\label{blockequal}
Under these hypotheses, $\tilde M=M$ is impossible.
\end{lem}
\begin{proof}
The blocks $A,D$ would have equal opposite moduli and could be made
real. A negative product of their entries would cancel their column
inner product, forcing respectively $j(C)=0$ or $j(B)=0$.
Make both blocks positive by independent phases; then $\eta=\pm1$.
For squared moduli of a block set
$f(X)=\sqrt{X_{11}X_{22}}-\sqrt{X_{12}X_{21}}$.
Differentiate $\det D=\eta\overline{\det A}$ along an arbitrary
unitary variation. Since $\eta'$ is imaginary, real parts give
\[
 df_{M_D}(\dot M_D)-\eta\,df_{M_A}(\dot M_A)=0
\]
This functional is nonzero on the 9-dimensional bistochastic tangent
space: its value on the checkerboard variation in $A$ is $-2\eta$.
Its kernel is therefore exactly $im(dp_U)$.
On blocks with equal opposite entries, the segment $M\to J_4$ scales
both $f$ values by $1-t$. Thus $J_4-M$ belongs to this kernel,
contradicting $dh_M(J_4-M)>0$.
\end{proof}
\begin{proof}[Proof of Theorem~\ref{threecase}]
Lemmas~\ref{onepair} and~\ref{twosingle} settle 1 or 2 lines.
For 3 lines place the repeated pair in $A$, so $j(A)=0$ and
$j(C)\ne0$. If $j(B)=0$, the first 2 rows use at most 2 lines.
Otherwise Lemmas~\ref{blockcurve} and~\ref{blockequal} give positive
second derivative. Transposition treats row pairs.
\end{proof}

The block direction can mix all 4 columns; existence of an increasing
single pair at every nonflat critical point remains open. For a 3-line column pair with
singleton rows $a,b$, orthogonality and stationarity give
$p_a^{-1}-q_a^{-1}=p_b^{-1}-q_b^{-1}$: project both equations onto a
linear functional killing the repeated line.

If every column pair uses at most 2 lines, any local maximum is flat
by Theorem~\ref{threecase}. As a separate finite check,
\texttt{phase\_grid.py} tests all $4^9=262144$ dephased fourth-root
phase patterns: $192$ pass the necessary sign-balance test, comprising
$168$ real non-Hadamard patterns, $6$ real and $18$ nonreal Hadamard
patterns. Passing this test does not prove unitary realizability.

\subsection{The equations for 4 distinct phase lines}\label{fourlinessection}
We now assume that all $36$ quantities $j(E)$ from
Theorem~\ref{threecase} are nonzero. The first-order equations already
restrict the moduli in this case. Write $e=(1,1,1,1)^t$ and put
\[
 R_{ia}=\frac1{|U_{ia}|},\qquad P=I-\frac14ee^t,
 \qquad G=PRP
\]
Thus $G$ is obtained from the reciprocal moduli by subtracting the row
and column means and adding the overall mean.

\begin{lem}\label{distinctdifferences}
At a nonflat critical point under our assumption, the $4$ numbers
$R_{ia}-R_{ib}$ are pairwise distinct for every pair of columns $a,b$.
The same holds for rows.
\end{lem}
\begin{proof}
Put $z_i=\bar U_{ia}U_{ib}$ and $d_i=R_{ia}-R_{ib}$. Orthogonality and
criticality give, respectively,
\[
 \sum_i z_i=0,\qquad \sum_i d_i z_i=0
\]
Indeed, the second equation is the $(a,b)$ entry of $S^*U-U^*S=0$.
If $d_1=d_2$, subtract $d_1$ times the first equation from the second.
The vectors $z_3,z_4$ are linearly independent over $\mathbb R$, so all
4 differences agree. Since the 2 columns have equal squared norms
and $r\to r^{-2}$ is strictly decreasing, their common difference
is zero. In particular, for every 2 rows $i,j$, the row difference
$R_{il}-R_{jl}$ agrees at $l=a,b$. Apply the same argument to
$U^t$, which is also critical. All row differences are zero, and column
normalization now gives $|U_{ia}|=1/2$ everywhere.
\end{proof}

The distinct reciprocal differences determine 2-dimensional planes
for each column pair. Compatibility between these planes gives a rank
restriction on the whole matrix.

\begin{thm}\label{reciprocalrank}
At a nonflat critical point of $L$ on $U_4$ with $4$ distinct phase
lines for every pair, the centered reciprocal modulus matrix $G$ has
rank exactly $2$.
\end{thm}
\begin{proof}
Fix $a,b$ and put $J_{kl,ab}=Im(\bar z_k z_l)$.
The 2 columns $Re(z),Im(z)$ span the
orthogonal complement in $\mathbb R^4$ of $e,d$. Both planes have
dimension 2 by our phase assumption and Lemma~\ref{distinctdifferences}.
The oriented area forms of complementary 2-planes are Hodge dual up
to a nonzero scalar. In coordinates their $2\times2$ minors therefore
satisfy
\[
 J_{kl,ab}=c_{ab}\,\epsilon_{ijkl}(d_i-d_j)
\]
where $i,j,k,l$ are distinct and $\epsilon_{ijkl}$ is the sign of this
permutation. The real number $c_{ab}$ is independent of the row indices.
Exchanging $a,b$ negates both $J$ and $d$, hence $c_{ba}=c_{ab}$.

For fixed $k,l,a$ we have $\sum_{b\ne a}J_{kl,ab}=0$, by orthogonality of
rows $k,l$. For the complementary rows $i,j$, put $v_a=R_{ia}-R_{ja}$.
The preceding identity gives
\[
 \sum_{b\ne a} c_{ab}(v_a-v_b)=0
\]
Let $H$ have off-diagonal entries $-c_{ab}$ and diagonal entries
$\sum_{b\ne a}c_{ab}$. This nonzero real symmetric matrix kills
$e$ and every row difference of $R$, viewed as a column vector.
It therefore kills their projections by $P$, which span the row space
of $G$. This space is orthogonal to $e$, so $rank(G)\le2$.

Rank zero contradicts Lemma~\ref{distinctdifferences}. In rank 1 we
could write
\[
 R_{ia}=u_i+v_a+x_i y_a
\]
with all parameters real, by restoring the row and column means to $G$.
The lemma forces the $x_i$ to be pairwise distinct, and likewise the $y_a$.
Criticality says that $iS$ is tangent at $U$, since $U^*S$ is self-adjoint.
Subtracting the phase tangents $i\,diag(u)U+iU\,diag(v)$ leaves
$i\,diag(x)U\,diag(y)$ tangent. Thus $U^*\,diag(x)U$ commutes with
$diag(y)$ and must be diagonal. Every column of $U$ would be an
eigenvector of $diag(x)$. Its eigenvalues are simple, whereas every
column has 4 nonzero entries.
\end{proof}

Interpolation at the 4 distinct reciprocal differences gives a
polynomial identity involving only the moduli.

\begin{prop}\label{pairequation}
Write $p_i=|U_{ia}|$, $q_i=|U_{ib}|$, $d_i=p_i^{-1}-q_i^{-1}$ and
$D_i=\prod_{j\ne i}(d_i-d_j)$. Under the assumptions of
Theorem~\ref{reciprocalrank}, every column pair satisfies
\[
 \sum_{i=1}^4 p_i^2 q_i^2 D_i=0
\]
The $6$ row pairs satisfy the same identity.
\end{prop}
\begin{proof}
For $k=0,1,2$, interpolation of $t^k$ at the 4 distinct nodes $d_i$
gives $\sum_i d_i^k/D_i=0$ by comparison of cubic coefficients.
The solutions of $\sum z_i=\sum d_i z_i=0$ therefore have the form
\[
 z_i=(A+B d_i)/D_i,\qquad A,B\in\mathbb C
\]
Indeed these vectors form a 2-dimensional complex space.
Now $p_i^2q_i^2D_i^2=|A+B d_i|^2$ is the value of a quadratic at $d_i$.
The coefficient of degree 3 in its interpolating polynomial is
$\sum_i p_i^2q_i^2D_i$, giving the formula.
\end{proof}

The same interpolation also reconstructs the relative phases.
It leaves only finitely many conjugation choices to check.

\begin{prop}\label{phaseuniqueness}
In the $4$-line case, prescribed moduli determine at most $1$ dephased
nonflat critical unitary, up to complex conjugation. There is an explicit
reconstruction test.
\end{prop}
\begin{proof}
For each pair interpolate $p_i^2q_i^2D_i^2$ at $d_i$. By
Proposition~\ref{pairequation} this gives a quadratic
$q(t)=a_2t^2+a_1t+a_0$. 4 distinct phase lines require and, for that
pair, are equivalent to
\[
 a_2>0,\qquad 4a_0a_2>a_1^2
\]
Take $B=\sqrt{a_2}$ and
$A=a_1/(2B)+i\sqrt{a_0-a_1^2/(4a_2)}$ in $z_i=(A+Bd_i)/D_i$.
This gives all pair solutions up to common phase and conjugation.
Pair column 1 with each other column and make their first entries
positive. The 3 conjugation choices leave at most 4 candidates
up to simultaneous conjugation; test their remaining orthogonality and
stationarity equations.

In fact at most 1 candidate survives. Let $v$ span
$\ker G\cap e^\perp$. The symmetric matrix $H$ in the proof of
Theorem~\ref{reciprocalrank} has this line as its image, so
$c_{ab}=\lambda v_av_b$ with all factors nonzero. Choose the conjugate
representative with $\lambda>0$. For each $b>1$, the sign
\[
 sgn(J_{12,1b})
 =sgn\big(v_1v_b(R_{31}-R_{3b}-R_{41}+R_{4b})\big)
\]
fixes the conjugation choice. This proves uniqueness. In particular,
symmetric moduli force $U^t=U$: the alternative $U^t=\bar U$ would make
$U$ self-adjoint and give a real principal $2\times2$ phase product.
\end{proof}

\subsection{Common profiles and at most 4 moduli}\label{fewmoduli}
The rank and pair equations have concrete consequences. We first handle
2 group patterns, then show that a common multiset reduces to them.

\begin{thm}\label{groupmoduli}
Let $U\in U_4$ be a local maximum of $L$. Suppose that after row and
column permutations its moduli have the form $|U_{ia}|=f(a-i)$ on
$\mathbb Z_4$ or on $\mathbb Z_2^2$. Then $U$ is flat.
\end{thm}
\begin{proof}
If $U$ were nonflat, all $j(E)$ would be nonzero by
Theorem~\ref{threecase}. The reciprocal moduli have the same symmetry.
In the cyclic case write their first row as $(a,b,c,d)$.
On the orthogonal complement of constants the eigenvalues of the
circulant are $a-b+c-d$ and the conjugate pair
$a-c\pm i(b-d)$. Rank 2 forces $a+c=b+d$. Write
\[
 (a,b,c,d)=(m+x,m+y,m-x,m-y)
\]
Read 2 adjacent rows cyclically, with reciprocal moduli
$(a,b,c,d)$ and $(b,c,d,a)$. Their 4 differences are
$(x-y,x+y,y-x,-x-y)$. Substitution in
Proposition~\ref{pairequation} and multiplication by $a^2b^2c^2d^2$
gives
\[
 64m^3xy(x-y)(x+y)=0
\]
For example, the left side is the expansion of
$c^2d^2D_1+d^2a^2D_2+a^2b^2D_3+b^2c^2D_4$. Since $m>0$, one of the
4 other factors must vanish. If $x=y$ or $x=-y$, 2 adjacent-row
differences coincide. If $x=0$ or $y=0$, 2 opposite-row differences
in $(2x,2y,-2x,-2y)$ coincide. Both conclusions contradict
Lemma~\ref{distinctdifferences}.

For $\mathbb Z_2^2$ the 3 nonconstant eigenvalues are
$a+b-c-d$, $a-b+c-d$, $a-b-c+d$. One must vanish. Relabel the 3
nonidentity elements to arrange $a+b=c+d$. The rows $(a,b,c,d)$ and
$(c,d,a,b)$ have differences $(a-c,b-d,c-a,d-b)$. The first and
last coincide, giving the same contradiction.
\end{proof}

A common row and column profile can be decomposed into permutation
matrices. At order 4 these recover exactly the 2 group patterns
just considered.

\begin{thm}\label{multisetmoduli}
Suppose that each row and each column of a local maximizer on $U_4$
contains the same multiset of entry moduli. Then the unitary is flat.
\end{thm}
\begin{proof}
For a modulus occurring $k$ times per row and column, join row $i$ to
column $a$ when the entry has this modulus. The resulting bipartite
graph is $k$-regular. Any set of rows has at least as many neighbors,
by counting incident edges. Hall's theorem gives a perfect matching,
whose removal lowers every degree by 1. Repeating this for each
modulus partitions the array into 4 permutation matrices, each
carrying a constant modulus. This step allows repetitions among the
4 moduli.

Permute columns to make the first matching the identity. The other
3 permutations have no fixed points. If they are all involutions,
they are the 3 double transpositions, which together with the
identity act regularly as the Klein 4-group. Otherwise one is a
4-cycle $s$. At each row $i$ the remaining columns are $s^2(i)$ and
$s^3(i)$. These edges form a single bipartite cycle of length 8. Indeed,
traversing an $s^3$ edge and returning along an $s^2$ edge advances
the row by $s$. Its only perfect matchings are the 2 alternating
ones. Thus the 4 permutations are $1,s,s^2,s^3$. In either case
we can label rows and columns by the group so the modulus at $(i,a)$
depends only on $a-i$. Theorem~\ref{groupmoduli} applies.
\end{proof}

In particular, a local maximizer on $U_4$ with at most 2 distinct entry
moduli is flat. Indeed, if these moduli are $r\ne s$, the number $k$ of
entries of modulus $r$ on any row or column satisfies
$kr^2+(4-k)s^2=1$. This fixes $k$ throughout the matrix.

With 3 distinct moduli, different row profiles can still occur.
Normalization and the reciprocal-difference restrictions rule them out.

\begin{prop}\label{threemoduli}
A local maximum on $U_4$ having at most $3$ distinct entry moduli is
flat.
\end{prop}
\begin{proof}
It remains to consider 3 distinct squared moduli, which we call
symbols. Suppose that a nonflat local maximum exists. By
Theorem~\ref{threecase} and Lemma~\ref{distinctdifferences}, 2 rows
can have equal symbols in at most 1 column. Also, wherever one row has a repeated symbol,
another row must take different values at those positions. Both rules
hold for columns.

A constant row is impossible, since another row would need 4
distinct symbols. Suppose that one row has the form $(x,x,x,y)$.
Each other row contains all 3 symbols in its first 3 positions.
Normalization forces their fourth entries to have the same squared
modulus $z$. Comparing the first row with the last column gives
$3x+y=3z+y$, hence $x=z$. Comparing with any other row now gives
$3x+y=x+y+w+x$, where $w$ is the third symbol. This forces $w=x$.
We have ruled out triples in both rows and columns.

The possible multiplicities in a row are now $(2,1,1)$ and $(2,2)$.
2 rows of the second kind must use the same pair $x,y$, since their
sums agree and there are only 3 symbols. Under a row
$(x,x,y,y)$ the second such row must put distinct entries in the first
pair and in the last pair. The rows then agree twice. Thus at most 1
row has type $(2,2)$, and likewise for columns. All rows and columns
of type $(2,1,1)$ repeat the same symbol, by normalization.

If there is a row $(x,x,y,y)$, comparison with a row of the other type
shows that its repeated symbol is the third symbol $z$ and that
$x+y=2z$. There are 6 copies of $z$ in the matrix. Consequently
exactly 1 column has type $(2,2)$ and contains no $z$. Delete this
column and the row $(x,x,y,y)$. Each row and column of the remaining
$3\times3$ array contains 2 copies of $z$. Any 2 of its rows
share a $z$ in 1 column. In the deleted column 2 of these 3
rows must also agree, since only $x,y$ occur there. This contradicts
the first rule. Transposition rules out a column of type $(2,2)$ as
well. All rows and columns therefore have the same multiset, and
Theorem~\ref{multisetmoduli} finishes the proof.
\end{proof}

4 distinct moduli require a finite classification. The count
constraints reduce the problem to $27$ patterns, each excluded by
explicit polynomial identities.

\begin{thm}\label{fourmoduli}
Every local maximum of $L$ on $U_4$ with at most $4$ distinct entry
moduli is flat.
\end{thm}
\begin{proof}
Only 4 distinct moduli and unequal profiles remain. Label symbols
$0,1,2,3$ in increasing reciprocal order, with squared moduli
$s_0>s_1>s_2>s_3$. For a sorted row or column profile $w$, put
$n_k(w)=\#\{w_i\ge k\}$ and $u=(s_0-s_1,s_1-s_2,s_2-s_3)>0$.
Its squared norm is $4s_0-n(w)\cdot u$, so all 8 profiles have the
same scalar product with $u$.

For enumerating profiles, $u$ can be taken integral with entries at
most $32$. In rank 2, cross 2 independent differences
$n(w)-n(w_0)$: the coordinates are bounded by $2\cdot4^2$ and have
one sign. In rank 1 a difference $d$ has both signs; assign the sum
of negative absolute entries to its positive positions, the sum of
positive entries to its negative positions, and $1$ to zeros.
Rank zero is the common-profile case. This imposes no arithmetic
restriction on the actual moduli.

There are $35$ profiles. Group them by their scalar product for
$u\in\{1,\ldots,32\}^3$, retaining groups with at least 2 profiles
and all 4 symbols: $157$ groups remain, $93$ maximal. Assemble
4 words from each maximal group, requiring columns from the same
word set. Reject row pairs repeating an ordered symbol pair or agreeing
twice, by Lemma~\ref{distinctdifferences}; reject fewer than 4
symbols or a common profile. Canonicalize by sorting rows after every
column permutation, also after transposition. Exactly $27$ arrays remain.

Scale the reciprocal levels to $1<x<y<z$. Impose equal row and column
sums of inverse squares, the 12 pair equations, and
\[
 \det(R_{ia}-R_{1a}-R_{i1}+R_{11})_{2\le i,a\le4}=0
\]
Clear positive denominators and factor over the rationals. A factor
with nonzero coefficients of one sign after
$x=1+a,y=1+a+b,z=1+a+b+c$ is nonvanishing for $a,b,c>0$.
This check excludes $10$ arrays by norms or determinant and $13$ by a
pair equation. In 3 more, the surviving determinant factors are
$3x-y-z-1$, $6x-4y-z-1$, $6x-y-4z-1$; each also requires
$x^2(1+y+z)-yz=0$. Substitution for $x$ makes this strictly positive
by the same coefficient test.

The last array and determinant equation are
\[
 \begin{pmatrix}0&2&3&3\\2&1&2&3\\3&2&1&2\\3&3&2&0\end{pmatrix},
 \qquad xz-x+y^2-3yz+y+z^2=0
\]
Rescale the reciprocal levels to $r_3=1$ and put
$h=1-r_0$, $t=(1-r_2)/h$. The determinant gives
\[
 (r_0,r_1,r_2,r_3)=(1-h,1-h(t+t^2),1-ht,1),
 \qquad 0<h,t<1,\quad t+t^2<1
\]
Take the norm difference of rows 1 and 2 and the pair equations for
rows $(1,2),(1,3)$. Discard sign-definite factors, homogenize with
$r_0$, and substitute these levels. Removing constant factors and
powers of $h,t-1,t+1$ leaves $F,G,K$ of degrees $5,6,6$ in $h$.
Their resultants have monic greatest common divisor
\[
 \gcd(Res_h(F,G),Res_h(F,K))
 =t^{26}(t^2+t-1)^4
\]
It cannot vanish in the required domain. This exhausts the $27$ arrays.
Every count, factor test and resultant is reproduced by
\texttt{four\_moduli.py}; \texttt{four\_lines.py} checks the earlier
cyclic factorization and matching classification.
\end{proof}

\section{Order 4: curvature and numerical evidence}\label{fourcurvature}
The restrictions above leave the case of 4 distinct phase lines and
at least 5 entry moduli. We examine its remaining curvature, then
give exact calibration saddles and the results of the searches.

\subsection{Curvature in the remaining case}\label{curvaturesection}
At a nonflat 4-line critical point $dp_U$ has rank 8, by
Lemma~\ref{eightblocks}. Its kernel consists of 7 phase directions
and 1 further direction, represented by $iS=iUX$: this fixes all
moduli to first order and is not a phase tangent, since that would give
$R_{ij}=u_i+v_j$ and $G=0$.
For an anchor $(a,b)$ set
\[
 T^{ab}_{ij}=R_{ij}-R_{ib}-R_{aj}+R_{ab},\qquad
 m_k^{ab}=\sum_{ij}|U_{ij}|^2(T^{ab}_{ij})^k
\]
The following formula tests this direction, even when phase lines coincide.

\begin{prop}\label{realcubic}
At every full-support critical point on $U_4$, all $16$ values
\[
 C_{ab}=m_3^{ab}-\frac32m_1^{ab}m_2^{ab}
              +\frac12(m_1^{ab})^3
\]
agree and equal $Q_U(X)=\sum_{ij}|U_{ij}|^{-1}-Tr(X^3)$, where
$X=U^*S$. A positive common value excludes a local maximum.
\end{prop}
\begin{proof}
Put $\alpha_i=R_{ib}$, $\beta_j=R_{aj}-R_{ab}$ and
$B=X-U^*\,diag(\alpha)U-diag(\beta)$. Then
$(UB)_{ij}=T^{ab}_{ij}U_{ij}$ and phase invariance gives $Q_U(B)=Q_U(X)$.
The self-adjoint matrix $B$ kills the independent vectors $e_b,U^*e_a$,
so it has rank at most 2. Thus
\[
 Tr(B)=m_1,\quad Tr(B^2)=m_2,\quad
 Tr(B^3)=\frac32m_1m_2-\frac12m_1^3
\]
Substitute $X=B+U^*\,diag(\alpha)U+diag(\beta)$ in the second variation.
The additive terms cancel against the corresponding part of
$\sum_{ij}|U_{ij}|T_{ij}^2$, leaving $m_3-Tr(B^3)$.
\end{proof}
In the 4-line case $B$ has rank exactly 2: rank 1 would make
the phase products in its nonzero complementary $3\times3$ block
real. Hence $m_2\ne m_1^2$ and $m_2\ge m_1^2/2$. The equalities
$C_{ab}=C_{11}$ give 15 necessary consistency equations, with no
claim of independence or sufficiency. On the $3596$ saved nonflat returns,
\texttt{curvature.py} finds discrepancy below $3.3\cdot10^{-11}$ from
the trace formula and common values in $(3.5867,3.9118)$.

\begin{prop}\label{rankonlyobstruction}
Bistochasticity, centered reciprocal rank $2$ and nonzero mixed
differences do not force a positive cubic: there are positive data $M$
with all $C_{ab}\in(-0.008,-0.000017)$. These data have no critical
unitary realization.
\end{prop}
\begin{proof}
An explicit certificate uses $G=10^{-10}AD^t$, where
\[
 A=\begin{pmatrix}-15233&-48634\\7967&-41084\\13665&113540\\-6399&-23822\end{pmatrix},
 \qquad
 D=\begin{pmatrix}79990&-3775\\-104383&15717\\37032&149\\-12639&-12091\end{pmatrix}
\]
Both factors are centered and have rank 2. Put
$R_{ij}=G_{ij}+u_i+v_j$, $v_4=0$, and $M_{ij}=R_{ij}^{-2}$.
Solve the 4 row and first 3 column normalization equations in
the box of radius $10^{-10}$ around
\[
 \begin{gathered}
 (u_1,u_2,u_3,u_4)=(2.007055173083,2.007781759308,\ 
                  2.008127664289,2.002764453315)\\
 (v_1,v_2,v_3)=(-0.001933545200,-0.000421319991,-0.004970340424)
 \end{gathered}
\]
Let $F$ be 1 minus these 7 sums. Take the inverse of $DF$ at the center and round its entries to 40
significant decimals, obtaining a rational preconditioner $Y$. Interval arithmetic verifies
$\|I-YDF\|_\infty<4.5\cdot10^{-9}$ and that $x-YF(x)$ maps the box
into its concentric box of radius $5\cdot10^{-13}$. The contraction
principle gives a unique solution. On the whole box $R>0$, all mixed
differences are nonzero, and the claimed cubic bounds hold.
Moreover $C_{11}\in(-0.002962,-0.002961)$ and
$C_{44}\in(-0.000017406,-0.000017404)$; their inequality excludes a
critical unitary by Proposition~\ref{realcubic}.
All interval operations are in \texttt{curvature.py}.
\end{proof}

There is nevertheless a positive result near $J_4$. Write
\[
 H=\begin{pmatrix}1&1&1\\1&-1&-1\\-1&1&-1\\-1&-1&1\end{pmatrix},
 \qquad F=(e\ H),\qquad G=HAH^t
\]
The 12 pair equations mean the equations of
Proposition~\ref{pairequation} with weights $M_{ia}M_{ib}$; they make
sense even when differences coincide.
\begin{thm}\label{regularflatlimit}
Let $M(t)$ be a $C^2$ path of positive bistochastic matrices with
$M(0)=J_4$, $M'(0)\ne0$. Suppose for all small $t>0$ that
$rank(G(t))\le2$, the $12$ pair equations hold, and all
$16$ $C_{ab}(t)$ agree. Write
$G(t)=tHA_0H^t+t^2HA_1H^t+o(t^2)$.
After row and column permutations and possibly transposition,
\[
 A_0=\begin{pmatrix}a&b&0\\0&0&0\\0&0&0\end{pmatrix},\quad
 a^2+b^2>0,\qquad C(t)=32(a^4+b^4)t^4+o(t^4)>0
\]
No unitary realization is assumed.
\end{thm}
\begin{proof}
Normalization gives $R=2ee^t+tG_0+O(t^2)$ and
$M=J_4-tG_0/4+O(t^2)$, where $G_k=HA_kH^t$.
For each anchor let $T,Z$ be the mixed differences of $G_0,G_1$.
Put $u_k=\frac14\sum T_{ij}^k$. The cubic coefficient is
$E=u_3-\frac32u_1u_2+\frac12u_1^3$.
For the matrix of these coefficients, direct expansion gives
\[
 (F^tEF/16)_{k0}=-96\prod_j(A_0)_{kj},\qquad
 (F^tEF/16)_{0k}=-96\prod_i(A_0)_{ik}\quad(k=1,2,3)
\]
They vanish, so every row and column of $A_0$ contains a zero.
The leading coefficient of a column-pair equation is
$32\prod_i(A_0v)_i$, with $v=e_j\pm e_k$ up to sign.
Thus any 2 columns $u,v$ satisfy
$\prod_i(u_i+v_i)=\prod_i(u_i-v_i)=0$.
They have a common zero: otherwise 2 different positions have
$u_i=\pm v_i\ne0$, forcing both third entries to vanish because each
column contains a zero. Their zero sets intersect pairwise. Unless
all 3 share a zero, they are the 3 distinct 2-element
subsets, making $A_0$ a nonsingular monomial matrix. This contradicts
$\det A_0=0$. Transposition gives a zero column as well.

For the remaining block $\left(\begin{smallmatrix}a&b\\c&d\end{smallmatrix}\right)$,
the $(1,1),(1,2),(2,1),(2,2)$ transformed cubic coefficients are
$48bcd,48acd,48abd,48abc$. At most 2 entries survive, and then $E=0$.
For completeness the fourth coefficient can be recovered without any
asymptotic guess: expand the same cubic expression using
\[
 T(t)=tT+t^2Z,\qquad M(t)=J_4-tG_0/4
\]
and retain degree 4. Call the resulting anchor matrix $D$.
If the 2 entries occupy different rows and columns, put
$A_0=diag(a,b,0)$. Then
\[
 (F^tDF/16)_{33}=-48ab(a^2+b^2)\ne0
\]
independently of $A_1$, contradicting equality of the anchor values.
The remaining support is 1 row or column. For the displayed $A_0$
and arbitrary $A_1$, the constant coefficient is
\[
 (F^tDF/16)_{00}=32(a^4+b^4)
\]
These polynomial identities are checked exactly in \texttt{flat\_limit.py}.
The $C^2$ remainders contribute $o(t^4)$, giving the conclusion.
\end{proof}
This result concerns regular paths in the modulus equations. It gives
neither a uniform punctured neighborhood nor a global exclusion.
Paths with vanishing first derivative can require fractional powers
after reparametrization and are not covered.

A bounded exploratory search parametrized $G=AB^t$ with
$\|G\|_F<\sqrt3$, balancing $R$ by additive row and column terms.
The $300$ starts comprised $30$ fitting cubic consistency alone, $30$
fitting pair equations alone, and $120$ each combining cubic consistency
with a negative-value target or with the pair equations. The last 2
minimum scaled residuals were $0.04015$ and $0.001853$, neither below
$10^{-7}$. These are auxiliary residuals, not unitary criticality tests.
\texttt{moduli\_search.py --verify-saved} rechecks every retained endpoint;
these searches exclude no additional critical unitaries.

It remains open whether normalization, reciprocal rank 2, the pair
identities, cubic consistency and phase compatibility force positive
common curvature at every nonflat 4-line critical point. A negative
common cubic alone would not settle the full Hessian test.

\subsection{2 exact calibration saddles}\label{fourcalibration}
2 of the stationary levels in the next subsection have explicit representatives.
For the first, take $a=\cos(\pi/8)/\sqrt2$ and
$b=\sin(\pi/8)/\sqrt2$. An exact representative is
\[
 U=\begin{pmatrix}
 a&b&a&b\\ a&-b&-b&-a\\ b&a&-a&b\\ b&-a&-b&a
 \end{pmatrix}
\]
The identities $2a^2+2b^2=1$ and $a^2-2ab-b^2=0$ give $UU^t=I$.
Multiplication gives $S^tU$ symmetric, proving criticality. We have
$L(U)=8(a+b)=4\sqrt{2+\sqrt2}$. The imaginary column rotation in
Lemma~\ref{realup} increases $L$. Along a real rotation of the same columns the
signs are locally constant, so the pair contributes $C\cos t+D\sin t$
with $C>0$ and $D=0$ by criticality. This gives a decrease.

For the highest observed nonflat level, use a bordered circulant matrix.
For $0<x<1$ write $s=\sqrt{1-x^2}$ and take the corner entry $x$, with
3 equal entries $s/\sqrt3$ along both borders. Take the inner block
circulant, with successive right shifts of the row
\[
 c_j(x,t)=\frac{-x+2i\sin(t+2\pi j/3)}{3},\qquad j=0,1,2
\]
The border and constant mode carry the block
$\left(\begin{smallmatrix}x&s\\s&-x\end{smallmatrix}\right)$.
The 2 nonconstant Fourier eigenvalues have modulus 1, since the sine
expands into the modes $e^{it}$ and $-e^{-it}$. This proves unitarity.
At $t=0$ the norm is
\[
 F(x)=2x+2\sqrt3\sqrt{1-x^2}+2\sqrt{x^2+3}
\]
We have $F''(x)=-2\sqrt3/(1-x^2)^{3/2}+6/(x^2+3)^{3/2}<0$.
The derivative is positive at zero and tends to $-\infty$ at 1, giving
a unique zero $x_0=0.6096097890\ldots$ and $F(x_0)=7.6376151868\ldots$.

To verify stationarity on all of $U_4$, both $S$ and $U$ at $t=0$
are self-adjoint and diagonal on the 2 nonconstant modes. On the
border and constant mode $S$ has diagonal entries $1,v$ and off-diagonal
entries $\sqrt3$, where $v=-1-2x/\sqrt{x^2+3}$. The 2 off-diagonal
entries of $S^*U$ differ by $s(1-v)-2\sqrt3 x=sF'(x)$, which vanishes
at $x_0$. Thus $S^*U$ is self-adjoint.

The $x$-direction decreases the norm. In the $t$-direction the varying
part of the norm is $\sum_{j=0}^2\sqrt{x^2+4\sin^2(t+2\pi j/3)}$.
Putting $r=\sqrt{x^2+3}$, its second derivative at zero is
\[
 \frac4x-\frac4r-\frac6{r^3}
 =\frac{6(r-x)(2r+x)}{xr^3(r+x)}>0
\]
This proves the saddle assertion independently of numerical Hessian
diagonalization. The 2 intermediate levels in the numerical table remain uncertified as exact critical values.
The 2 representatives are matched by \texttt{models.py};
\texttt{three\_lines.py} also checks the block variation and certifies
column curvature $M_{23}\in(0.6092,0.6094)$,
row curvature $M_{14}\in(-0.6677,-0.6674)$ and block-curve curvature
in $(1.4478,1.4480)$ on partition $12\mid34$ for this calibration. Thus the singleton rows
of a 3-line pair need not themselves give an increase.

\subsection{Numerical evidence and its limits}\label{foursearch}
A useful bound separates potential local maxima from the nonsmooth
boundary. We use Frobenius-unit generators as in Lemma~\ref{twocol}.

\begin{prop}\label{entrybound}
At a local maximum on $U_N$, $N\ge2$, every entry has modulus at least
\[
 b_N=\sqrt{N(N-1)^2+1}-(N-1)\sqrt N
\]
In particular, $b_4=\sqrt{37}-6$ and $b_5=9-4\sqrt5$.
\end{prop}
\begin{proof}
For an entry of modulus $p$, choose a partner of modulus $q$ in its row
with $q^2\ge(1-p^2)/(N-1)$. Choose the phase of the generator in Lemma~\ref{twocol} to agree with
$h_i$ in the selected row. Its contribution to the first
term of the curvature is then $(q^2/p+p^2/q)/2$. Each column has
$1$-norm at most $\sqrt N$, giving
\[
 M_{jk}\ge\frac{q^2}{2p}-\sqrt N
 \ge\frac{1-p^2}{2(N-1)p}-\sqrt N
\]
At a local maximum $M_{jk}\le0$. Solving the resulting quadratic
inequality gives the bound.
\end{proof}

We use the diagnostics
\[
 g(U)=\frac12\|S^*U-U^*S\|_F,\qquad
 d(U)=\max_{ij}\big||U_{ij}|-1/2\big|
\]
2 Apple M3 laptops (8 GB, 3 workers each)
ran 7-hour campaigns, with seeds $421701,421702$. Ascent uses Haar
starts, geodesic gradient steps and Armijo backtracking, capped at
$1800$ iterations. Stationary search uses damped least squares for
$(S^*U-U^*S)/(2i)=0$. 4 protocols start from Haar, phase-polar,
near-real and near-flat matrices; 4 first steer Haar matrices towards
$L\le c$, $c\in\{7.40,7.65,7.90,7.995\}$, with at most $24$ steering
and $120$ root steps.

The tolerances are $d\le10^{-6}$ for flatness, $g\le10^{-12}$ for
stationarity, and $10^{-8}$ for near-zero entries. At accepted nonflat
returns both full and 9-dimensional horizontal Hessians are checked;
top horizontal eigenvalue at most $10^{-7}$ triggers a candidate flag.
\begin{center}
\begin{tabular}{lrr}
\toprule
Endpoint status & Ascent & Stationary\\
\midrule
Numerically flat &83\,009\,628&19\,409\,961\\
Nonflat with positive curvature &0&26\,672\,800\\
Near zero &0&6\,126\,915\\
Unresolved &67\,268&27\,748\\
\midrule
Total &83\,076\,896&52\,237\,424\\
\bottomrule
\end{tabular}
\end{center}
There were no candidate, unlisted-level, numerical-error or
unitarity-error flags. All nonflat stationary returns had $L$ within
$2\cdot10^{-6}$ of one of the following 4 values. Repeated visits
to equivalent matrices are included in the counts.
\begin{center}
\begin{tabular}{rrr}
\toprule
$L$ (rounded) & Returns & Minimum top horizontal eigenvalue\\
\midrule
7.391036 &6\,764\,846 &2.2961006\\
7.633239 &9\,497\,039 &2.8149691\\
7.633377 &7\,875\,896 &2.6601070\\
7.637615 &2\,535\,019 &2.8768867\\
\bottomrule
\end{tabular}
\end{center}
Every unresolved ascent endpoint satisfies
$d\le6.01386943033244\cdot10^{-6}$; not all were continued, so this
does not assert convergence of every trajectory. For stationary returns,
\[
 \max_{\text{unresolved }U}\min_{ij}|U_{ij}|
 =5.6974893208468939\cdot10^{-7}
\]
All unresolved and near-zero stationary endpoints thus fall far below
$\sqrt{37}-6=0.08276\ldots$. This comparison has a large numerical
margin but is not an interval certificate for the entire campaign.

The initial $8192$-start survey gave $3596$ nonflat uphill, $4002$
near-flat, $589$ near-zero and $5$ unresolved returns. The 4 nonflat
groups have numerical positive indices $3,3,2,1$. For every nonflat
return, interval arithmetic proves a recorded pair curvature exceeds
$0.60$ throughout its coordinate box of radius $10^{-8}$. This excludes
local maxima in the boxes without proving existence or uniqueness of
critical points there. The other $594$ endpoints have a smallest entry
at most $1.18\cdot10^{-8}$. A separate $1024$-start polar-ascent control
reached numerical flatness, including $11$ final least-squares refinements.

The compact archives retain all $95\,016$ unresolved overnight endpoints,
aggregate counts and extrema, and representative outputs. Independent
audits recheck $99\,164$ retained endpoints and all $98\,914$ stored
spectrum pairs. The sources are \texttt{search.py} and \texttt{audit.py} for the initial
survey and its certificates, and \texttt{ahc\_ascent.py} and
\texttt{ahc\_stationary.py} for the campaigns. \texttt{compact.py}
exports and verifies their archives; \texttt{audit\_compact.py}
recomputes all retained diagnostics. The archives contain their exact
generating sources and complete aggregate counts, but ordinary discarded
endpoints are unavailable.
These searches sample basins, not stationary points uniformly; repeated
visits are not independent evidence of completeness. No covering of
$U_4$ is claimed, and an unvisited basin remains possible.

\section{Order 5: exact models and gradient geometry}\label{fivegeometry}
\subsection{A saddle invisible to pair tests}
The order-5 search needs the full Hessian. Its main calibration
shows why pair rotations alone cannot decide local maximality.
For $0<x<1$, put
\[
 a=(1+x)/4,\qquad b=\sqrt{1-x^2}/2,\qquad c=(1-x)/4
\]
Let $U(x)$ have corner entry $x$, 4 entries $b$ on each border, and
an inner circulant block whose first row is
\[
 (-a,\ c-i/2,\ -a,\ c+i/2)
\]
The other inner rows are successive right shifts. On the border and
normalized constant inner vector, $U(x)$ acts by
\[
 \begin{pmatrix}x&\sqrt{1-x^2}\\\sqrt{1-x^2}&-x\end{pmatrix}
\]
Its other Fourier eigenvalues are $1,-1,-1$, in an appropriate order.
Thus $U(x)$ is self-adjoint and unitary. Counting moduli gives
\[
 F(x)=L(U(x))=2+3x+4\sqrt{1-x^2}+2\sqrt{(1-x)^2+4}
\]
Full criticality reduces here to the scalar equation $F'(x)=0$.

\begin{prop}\label{fivecalibration}
There is a unique $x_0\in(0,1)$ for which $U_0=U(x_0)$ is critical on
$U_5$. At $U_0$ every $2$-row and $2$-column orbit has a strict local
maximum modulo phases, but a $3$-column variation increases $L$.
The full Hessian has inertia $(15,9,1)$.
\end{prop}
\begin{proof}
Both $U(x)$ and its phase matrix $S$ are self-adjoint and diagonal on the
nonconstant inner Fourier modes. On the remaining 2-dimensional
space write
\[
 S=\begin{pmatrix}1&2\\2&v\end{pmatrix},\qquad
 v=-2+\frac{2(1-x)}{\sqrt{(1-x)^2+4}}
\]
The difference of the off-diagonal entries of $S^*U$ is
$\sqrt{1-x^2}(1-v)-4x=\sqrt{1-x^2}F'(x)$. Hence full criticality is
equivalent to $F'(x)=0$. We have
\[
 F'(x)=3-\frac{4x}{\sqrt{1-x^2}}-
                 \frac{2(1-x)}{\sqrt{(1-x)^2+4}}
\]
Its derivative is
\[
 F''(x)=-\frac4{(1-x^2)^{3/2}}
             +\frac8{((1-x)^2+4)^{3/2}}<0
\]
Since $F'(0)>0$ and $F'(x)\to-\infty$ as $x\to1$, the root is unique.
It is $x_0=0.5374606822272723\ldots$, with $L(U_0)=11.0911152536\ldots$.

Index the border by $0$ and the inner coordinates by $1,2,3,4$.
Cyclic symmetry leaves 3 column-pair types; self-adjointness gives the
same row tests. Lemmas~\ref{variations} and~\ref{twocol}, evaluated with
outward rounding, give
\begin{center}
\begin{tabular}{lc}
\toprule
Direction & Curvature enclosure\\
\midrule
Border--inner pair & $-0.8502450<M_{01}<-0.8502447$\\
Adjacent inner pair & $-0.3694043<M_{12}<-0.3694040$\\
Opposite inner pair & $-0.1978585<M_{13}<-0.1978582$\\
3-column generator $B$ & $0.5674395<Q_{U_0}(B)<0.5674400$\\
\bottomrule
\end{tabular}
\end{center}
Here $B$ is self-adjoint with $B_{12}=B_{23}=i$, and all other upper and
diagonal entries zero. Its Frobenius norm is $2$. The pair bounds are
for unit generators. The negative pair bounds and positive 3-column
bound prove the local assertions.

For the full inertia take $E_{jk}+E_{kj}$ and $i(E_{jk}-E_{kj})$,
in that order, on the first 8 pairs $j<k$ in lexicographic order. Polarizing $Q$
gives a $16\times16$ symmetric matrix. Put
$l=0.53746068222727233271039478$. Endpoint signs enclose $x_0$ in
$[l,l+10^{-26}]$. Interval elimination by congruence in the indicated
order gives pivot $14$ in $(1.58,1.59)$ and the other 15 pivots in
$(-3.02,-0.07)$. These 16 nonzero directions, together with the
9 phase zero modes, exhaust the tangent space. The root signs,
curvature bounds and elimination are reproduced by \texttt{code011/check.py}.
\end{proof}

\subsection{Further calibrations and bordered circulants}
The 4 original decimal seeds in \texttt{check.py} give the following
numerical refinements; inertias use zero threshold $10^{-6}$.
\begin{center}
\begin{tabular}{lrrrr}
\toprule
Seed & $L/(N\sqrt N)$ & Inertia & Top curvature & Best pair\\
\midrule
First $N=5$ & .99109 & $(14,9,2)$ & .32297 & $-.26412$\\
Second $N=5$ & .99202 & $(15,9,1)$ & .65312 & $-.19786$\\
$N=6$ & .98832 & $(22,11,3)$ & .65004 & $-.02548$\\
$N=7$ & .96837 & $(27,13,9)$ & 2.49105 & $+.24721$\\
\bottomrule
\end{tabular}
\end{center}
Polar projection and horizontal Newton refinement give residuals below
$10^{-14}$. Positive-curvature witnesses are separately interval-certified
on all 4 original coordinate boxes of radius $10^{-10}$.
For the second seed, row order $(5,1,2,4,3)$, column order $(2,1,3,4,5)$
and dephasing recover $U_0$ to $2\cdot10^{-15}$.
For the first, orders $(3,1,4,2,5)$ and $(2,3,5,4,1)$ give a symmetric
matrix invariant under simultaneous $(2\,3)(4\,5)$, to $3\cdot10^{-14}$.
If its positive eigenvalues are $u,v$ and negative ones $-a_j$, a
Frobenius-isotropic Gaussian direction has predicted uphill probability
$2.67\cdot10^{-6}$. It is the angular mean of
$\prod_j(1+a_j/d)^{-1/2}$, where $d=u\cos^2\theta+v\sin^2\theta$;
condition on the positive-plane radius, whose square is exponential of
mean 2. This explains why random direction tests can miss a saddle.

For $N=6$, column order $(4,2,6,5,1,3)$ and dephasing give symmetry to
$3\cdot10^{-15}$. 9 of 20 column triples have positive curvature;
the best is $.32062$ on original columns $(1,4,6)$. On that triple,
the self-adjoint generator with upper entries $(1-i,1+i,-1+i)$ has
$1.0604779<Q<1.0604782$ on the original rounding box.
For $N=7$, columns $(2,7)$ already give an increase. No exact
order-6 or order-7 critical model is asserted.

More generally, a border $x,\sqrt{1-x^2}/\sqrt m$ and circulant core
with Fourier eigenvalues $-x,z_1,\ldots,z_{m-1}$, $|z_k|=1$, give a
unitary of order $m+1$. At full support, $0<x<1$, stationarity in $x$
and all these phases is equivalent to full criticality: if the
phase eigenvalues of the core are $v_k$, phase stationarity makes
$\bar v_kz_k$ real for $k>0$; reality of $Tr(S^*U)$ then makes $v_0$
real, and the remaining skew entry is $\sqrt{1-x^2}\,\partial_x L$.
Restricted phase searches alone do not give this conclusion.

\subsection{The actual unstable branches}
Return to the exact order-5 model $U(x)$ from
Proposition~\ref{fivecalibration}. Replace its inner row by
\[
 (-a+\frac i2\sin t,\ c-\frac i2\cos t,
  -a-\frac i2\sin t,\ c+\frac i2\cos t)
\]
to obtain $U(x,t)$. Its nonconstant Fourier eigenvalues are
$e^{it},-1,-e^{-it}$. Put
$A=(a^2+\sin^2t/4)^{1/2}$, $C=(c^2+\cos^2t/4)^{1/2}$ and
$H(x,t)=x+4\sqrt{1-x^2}+8(A+C)$.
This surface is the local common fixed set of
$V\to PVP^{-1}$ and $V\to\bar V P^2$, where $P$ cycles the
4 inner coordinates. The first equation gives a circulant block;
the second makes the border block real and imposes the 3 displayed
Fourier phases. Both maps preserve $L$, so its gradient is tangent to
the surface. The induced metric and gradient equations are
\[
 \frac{2\,dx^2}{1-x^2}+2\,dt^2,\qquad
 \dot x=\frac{1-x^2}{2}H_x,\quad \dot t=\frac12H_t,
 \quad H_t=\sin(2t)(1/A-1/C)
\]
A fixed-corner $t$-rotation is a different curve: its corner stays
$x_0$, and it cannot become flat. Starting exactly at $U_0$ also gives
no motion. The proposition concerns the actual nonconstant gradient
trajectories with backward limit $U_0$.

\begin{prop}\label{branches}
The $2$ unstable branches through $U_0$ for gradient ascent converge to
$U(1/\sqrt5,\pm t_*)$, where $t_*=\frac12\arccos(1/\sqrt5)$.
Both limits are Fourier-equivalent.
\end{prop}
\begin{proof}
The restricted Hessian has opposite signs in $x,t$, and the full
horizontal Hessian has exactly 1 positive eigenvalue. The unstable
manifold theorem \cite{T}, on the phase quotient, therefore identifies
these 2 surface branches with the full unstable branches. Equivariance
and the trivial effective phase stabilizer fix their lifts through $U_0$.

Consider the positive branch in $0<x<1$, $0<t<\pi/2$.
On the boundary, $H\le F(x_0)$: the $t=0$ edge is $F(x)$, the
$t=\pi/2$ edge has a larger $|A^2-C^2|$ with the same sum, and
on $x=0,1$ use
\[
 H\le G(x):=x+4\sqrt{1-x^2}+4\sqrt{3+x^2},\quad
 G(0)=4+4\sqrt3,\quad G(1)=9<F(1/2)
\]
Here also $G(0)<F(1/2)=7/2+2\sqrt3+\sqrt{17}$.
The branch has strictly larger value than $F(x_0)$ and is trapped in
a compact interior superlevel set. Every forward limit point is
critical. Interior criticality gives $A=C$, hence $\cos(2t)=x$, and
then $G'(x)=0$. Since $G''<0$, the unique solution is $x=1/\sqrt5$.
Reflection treats the other branch. At the positive limit, the inner
row of $\sqrt5U$ is $(w^2,w^4,w^3,w)$ with $w=e^{2\pi i/5}$;
its cyclic shifts form the Fourier multiplication table, with row order
$(2,1,3,4)$ and column order $(1,2,4,3)$. The negative limit is its adjoint.
\end{proof}

\section{Order 5: numerical evidence and certification}\label{fivecampaign}
\subsection{The 7-hour campaigns}
The standalone \texttt{ahc5\_ascent.py} and
\texttt{ahc5\_stationary.py} check the full $25$-dimensional
Hessian and its $16$-dimensional horizontal restriction at accepted
nonflat stationary returns. They use flatness tolerance $10^{-6}$,
stationarity tolerance $10^{-10}$, near-zero threshold $10^{-7}$ and
positive-curvature threshold $10^{-7}$. Up to $80$ polishing steps
target residual $10^{-12}$. Ascent uses Haar, phase-polar,
Haar-polar, near-real, near-block and near-saddle starts, with at most
$1600$ steps. The stationary search uses Haar, phase-polar, near-real,
near-flat, near-saddle, near-block, bordered and circulant starts, and
4 steering caps $L/(5\sqrt5)\le .90,.95,.985,.999$, with at most
$32$ steering and $160$ root steps. The algorithms passed the derivative,
calibration and restart checks in \texttt{test\_ahc5.py}.
2 Apple M3 laptops (8 GB, 3 workers each)
completed 7-hour runs, with seeds $521811,521812$:
\begin{center}
\begin{tabular}{lrr}
\toprule
Endpoint status & Ascent & Stationary\\
\midrule
Numerically near flat &181\,164\,464&2\,402\,576\\
Nonflat with positive curvature &0&18\,997\,063\\
Near zero &0&1\,129\,156\\
Unresolved &0&404\,581\\
\midrule
Total &181\,164\,464&22\,933\,376\\
\bottomrule
\end{tabular}
\end{center}
There were no candidate, uncertain-Hessian or numerical-error flags.
Every ascent endpoint satisfied
$\max_{ij}\big||U_{ij}|-1/\sqrt5\big|<10^{-6}$.
The smallest recorded top curvature at a nonflat stationary return was
$0.206350429\ldots$; all such returns had residual below $10^{-12}$.
The earlier 5-minute local pilots gave $593\,024$ near-flat ascent
returns and $105\,248$ stationary returns: $10\,996$ near-flat,
$87\,250$ uphill, $5162$ near-zero and $1840$ unresolved. They are
separate runs; their $3840$-start test prefixes are not counted again.

One retained numerical critical point has $L=11.081082511058\ldots$,
top Hessian eigenvalue $0.206350429\ldots$, and largest row or column
pair curvature $-0.226683\ldots$. Its numerical inertia is again
$(15,9,1)$. An independent interval check proves
\[
 0.20635<Q_V(B)<0.20636
\]
for a fixed rational self-adjoint generator throughout the coordinate box
of radius $10^{-12}$ around the stored matrix. This excludes every
unitary in the box as a local maximum. Theorem~\ref{criticalcatalogue}
certifies an exact critical point near the representative of group $39$,
with this inertia and negative pair curvatures. The matrix and witness are in
\texttt{code011/data/ahc5-weak.json}, and the check is \texttt{code011/weak5.py}.

The near-zero returns lie below $b_5=9-4\sqrt5>0.0557$ and hence
cannot themselves be local maxima. The unresolved returns have
residual at least $1.254\cdot10^{-5}$, with smallest-entry moduli
reaching $0.1539657$; the entry bound does not dismiss all of them.
All $288$ retained unresolved endpoints have numerically positive
geodesic curvature, with smallest top value $3.67446\ldots$.
The saved-point follow-up below explains these stalls. A subsequent
partial replay addresses the original starting sequence; neither
calculation asserts exhaustive coverage.

The archives \texttt{code011/overnight5/ascent-mac.zip} and
\texttt{stationary-mac.zip} contain the generating sources, all-run
counters and extrema, and $96$ and $1076$ retained endpoints. Ordinary
endpoints use bounded sampling, supplemented by priority pools for weak
saddles and small unresolved residuals; every candidate flag would be
retained. The 2 databases total $4\,460\,544$ bytes.
\texttt{audit\_ahc5.py} rechecks source hashes, counters, starts and
retained endpoints, including $452$ stored pairs of full and horizontal
spectra and their uphill witnesses. Maximum spectral and witness
discrepancies were below $6\cdot10^{-14}$. The audit also computes
curvature at the retained unresolved points without refining them.
The unretained endpoints cannot be audited individually from these archives.

\subsection{Catalogue and systematic certificates}
The $452$ retained nonflat stationary returns form $40$ numerical groups
under row and column phases, permutations and complex conjugation, using
Frobenius tolerance $10^{-6}$. The comparison checks all pairs of row
and column permutations after dephasing; equal norms alone are not
used as equivalence tests. The table gives each representative norm,
its positive horizontal index $p$, and a lower bound $q$ for a fixed
uphill direction throughout every retained rounding box in that group.
Theorem~\ref{criticalcatalogue} certifies these indices and the displayed
norms at exact critical points near the $40$ representatives.
\begin{center}
\begin{tabular}{rrrr@{\qquad}rrrr}
\toprule
Group & $L$ & $p$ & $q$ & Group & $L$ & $p$ & $q$\\
\midrule
1 & 9.779960 & 6 & 17.279 & 21 & 10.475603 & 6 & 3.673\\
2 & 9.944272 & 6 & 3.572 & 22 & 10.475618 & 5 & 4.844\\
3 & 10.204351 & 7 & 5.864 & 23 & 10.475618 & 5 & 4.844\\
4 & 10.204611 & 6 & 4.582 & 24 & 10.475787 & 4 & 9.810\\
5 & 10.204611 & 6 & 4.582 & 25 & 10.493467 & 4 & 8.293\\
6 & 10.229297 & 6 & 6.638 & 26 & 10.599732 & 5 & 4.534\\
7 & 10.231276 & 5 & 7.045 & 27 & 10.656854 & 5 & 6.094\\
8 & 10.234788 & 5 & 3.557 & 28 & 10.658120 & 4 & 6.675\\
9 & 10.252396 & 6 & 10.519 & 29 & 10.659153 & 4 & 7.120\\
10 & 10.255519 & 6 & 3.761 & 30 & 10.662660 & 3 & 4.602\\
11 & 10.264114 & 6 & 3.507 & 31 & 10.662660 & 3 & 4.602\\
12 & 10.282845 & 6 & 4.288 & 32 & 11.000000 & 5 & 0.499\\
13 & 10.283393 & 5 & 3.930 & 33 & 11.047758 & 4 & 0.539\\
14 & 10.286257 & 5 & 6.297 & 34 & 11.048168 & 3 & 0.621\\
15 & 10.290855 & 5 & 5.992 & 35 & 11.048247 & 2 & 0.652\\
16 & 10.290855 & 5 & 2.989 & 36 & 11.061538 & 2 & 0.292\\
17 & 10.322914 & 4 & 36.252 & 37 & 11.079258 & 3 & 0.740\\
18 & 10.325913 & 4 & 6.172 & 38 & 11.080767 & 2 & 0.322\\
19 & 10.334241 & 4 & 8.302 & 39 & 11.081083 & 1 & 0.206\\
20 & 10.334241 & 4 & 8.302 & 40 & 11.091115 & 1 & 0.653\\
\bottomrule
\end{tabular}
\end{center}
For every one of the $452$ matrices, \texttt{catalogue5.py} rounds a
self-adjoint uphill generator to 12 decimals and interval-evaluates
$Q_U(B)$ on the real and imaginary coordinate box of radius $10^{-8}$.
Every lower bound exceeds $.20634$. These are rigorous exclusions of
local maxima among unitaries in the boxes. The grouping remains
numerical; these box exclusions alone do not establish critical-root
existence. That additional conclusion is proved below for the $40$
representatives. No completeness of the catalogue is asserted.
Representatives, matching
permutations and individual certificates are in \texttt{data/catalogue5.npz}
and \texttt{data/catalogue5.json}.

Allowing transposition merges groups $4/5$, $19/20$, $22/23$ and $30/31$
numerically, with matching errors below $6\cdot10^{-11}$, leaving $36$
groups. Groups $15,16$ have norms agreeing to $12$ decimals,
$L=10.290855136357$, but different smallest entry moduli,
$0.109049\ldots$ and $0.211227\ldots$. These comparisons are also
checked by \texttt{catalogue5.py}; exact equality of these 2 norms is not asserted.

\subsection{Stalled stationary searches}\label{fivestalls}
Write $R(U)=(S^*U-U^*S)/(2i)$ and $\Phi(U)=\|R(U)\|_F^2/2$.
Restarting the residual solver on the $288$ saved points, with damping
reset to $.01$ and up to $1600$ further steps, gives 3 outcomes:
\begin{center}
\begin{tabular}{rll}
\toprule
Count & Outcome & Diagnostic\\
\midrule
$139$ & Stationary saddles & $8$ at $L=10.475603$, $131$ at $L=10.475787$\\
$134$ & Positive residual minimum & $L=10.2696684693$, $\|R\|_F=0.0164775744$\\
$15$ & Near the zero-entry boundary & Smallest entry below $4.865\cdot10^{-7}$\\
\bottomrule
\end{tabular}
\end{center}
All $139$ in the first group had exhausted the original $160$ steps
and missed the polishing trigger $\|R\|_F<10^{-5}$. Restarting needs
only $7$--$29$ steps, giving residual below $10^{-10}$ and top curvature
above $3.67$. The $15$ boundary points already had an entry below
$4.269\cdot10^{-5}$ before restarting, far below $b_5$.
The $134$ middle returns agree with one matrix up to phases,
permutations and conjugation within $10^{-6}$. At each, replace $U$ by
\[
 U\exp\big(-0.6iR(U)/\|R(U)\|_F\big)
\]
and allow $800$ root steps. All reach the known level
$L=10.2641137511\ldots$, with top curvature above $3.50$ and residual
below $10^{-10}$. Perturbations of lengths $.1$ and $.3$ returned to
the residual minimum. The 3 final saddle
types match already retained representatives within $10^{-6}$.
These are numerical observations, reproduced by \texttt{stalled5.py},
which saves the matrices, equivalences and witnesses in
\texttt{data/stalled5.npz} and \texttt{data/stalled5.json}.
The identity $\nabla\Phi=Hess(L)\nabla L$ explains how
the residual can be stationary while $L$ is not. The positive minimum is certified in
Proposition~\ref{residualminimum}.

\subsection{Daytime follow-up and targeted search}\label{fivefollowup}
The program \texttt{ahc5\_day\_replay.py} replays seed $521812$ with
the original batches and solver. Unresolved full-support returns above
$b_5-10^{-6}$ receive $320$ restart steps, then if necessary the
length-$.6$ perturbation and up to $800$ root steps. Entry-bound
exclusions and nonstationary positive curvature are counted separately.
The complementary \texttt{ahc5\_day\_search.py}, seed $521819$, uses
8 protocols: root search near all $40$ representatives, near the
8 weakest, and along their soft directions; perturbed bordered and
circulant starts; geodesic and polar ascent near weak saddles; and
Haar root controls. Amplitudes range from $10^{-4}$ to $1$.
Both jobs check full and horizontal Hessians and retain new signatures.

Each Mac ran for 7 active hours with 2 workers and 1 BLAS thread per
worker. Unresolved identifiers and diagnostics and all candidate flags
are retained, with bounded pools for ordinary matrices. Data growth
stops at $112$ MiB, below the $128$ MiB cap. Baseline replay, standalone
launches, rollback, eigenbasis compatibility, locks and exact resumption
are checked by \texttt{test\_day5.py}. The $18$-second release checks
($5600$ replay and $10\,032$ targeted starts, no unresolved or candidate
returns) are prefixes, not additional starts. The completed runs give
the following counts; the lower block partitions the initial stalls
already included above.
\begin{center}
\begin{tabular}{lrr}
\toprule
Endpoint status & Replay & Targeted search\\
\midrule
Numerically near flat &1\,696\,013&10\,068\,469\\
Nonflat stationary, positive curvature &13\,672\,742&24\,303\,512\\
Small-entry exclusion &820\,253&579\,411\\
Unresolved or candidate flags &0&0\\
\midrule
Total starts &16\,189\,008&34\,951\,392\\
\midrule
Initially unresolved &284\,297&235\,667\\
Restart alone gives an uphill saddle &12\,306&15\,090\\
Escape and refinement give an uphill saddle &248\,455&146\,623\\
Excluded by the entry bound &23\,536&73\,954\\
\bottomrule
\end{tabular}
\end{center}
No uncertain-Hessian, unitarity or numerical-error flags occurred.
Replay covers $70.59\%$ of the original $22\,933\,376$ starts;
$6\,744\,368$ remain unvisited, so resolution of all $404\,581$
original stalls is not claimed. Its $284\,297$ stalls are those observed
on replay, and neither replay nor follow-up iterates add independent trials.
All $8\,737\,840$ targeted ascent returns are near flat. Every accepted
nonflat stationary return matches a stored signature, using sorted
moduli and horizontal spectra within $2\cdot10^{-5}$; both runs
encounter all $40$, with no new signature and least top curvature
$0.206350429\ldots$. Signature matching is screening, not an
equivalence proof or interval certificate.

The unchanged archives \texttt{code011/daytime/replay-mac.zip} and
\texttt{search-mac.zip} contain $788$ and $569$ retained runs, with
starting, raw and final matrices. Database sizes are $31.94,26.03$ MiB;
sums of parent and worker peak memory are $128.5,157.2$ MiB.
The program \texttt{audit\_day5.py} checks hashes, database integrity,
counters, retained matrices and $452+356$ full/horizontal spectrum
pairs and uphill witnesses. Polarization independently checks each
full Hessian; all spectral and witness errors are below $6\cdot10^{-14}$.
The adjacent \texttt{replay-mac-audit.json} and
\texttt{search-mac-audit.json} record these checks, which do not audit
unretained matrices.

All $788$ saved replay starts and $496$ targeted starts regenerate
within $10^{-11}$. The other $67$ soft-direction and $6$ ascent starts
depend on eigenbasis choices. Recover $B=-i\log(U_0^*U)$ from the saved
start $U$ and catalogue centre $U_0$. For a soft start of length $\rho$,
Gaussian coefficients $c_1,c_2,c_3$ and selected eigenvalues
$\lambda_1,\lambda_2,\lambda_3$, the audit checks
\[
 \|P_\lambda B\|_F^2
 =\rho^2\frac{\sum_{a:\lambda_a=\lambda}c_a^2}
                   {\sum_{a=1}^3c_a^2}
\]
Here $P_\lambda$ is the horizontal Hessian eigenspace projection;
eigenvalues within $10^{-7}$ are grouped. For ascent starts it checks
$B=\rho(W+.1E)/\|W+.1E\|_F$, where $E$ is the regenerated noise
and $W$ a unit vector in the top eigenspace. Relative norm and
projection discrepancies are below $2\cdot10^{-12}$. These are
numerical compatibility checks, not identical regeneration across
numerical libraries; saved starts and generating sources are unchanged.

\subsection{A positive minimum of the stationarity residual}
The following certificate explains the middle group in
Section~\ref{fivestalls}. It establishes an obstruction to residual
minimization, without producing a local maximum of $L$.

\begin{prop}\label{residualminimum}
On $U_5$, $\Phi$ has a strict local minimum modulo row and column phases
at a full-support point $V$ satisfying
\[
 .01647<\|R(V)\|_F<.01649
\]
In particular, this point is not stationary for $L$.
\end{prop}
\begin{proof}
The following integer matrix specifies a rational self-adjoint matrix $A$:
put $A_{jj}=10^{-12}C_{jj}$ and
$A_{jk}=10^{-12}(C_{jk}+iC_{kj})$ for $j<k$.
\[
\setlength{\arraycolsep}{3pt}
C=\begin{pmatrix}
495152000271&190465601450&-262733011052&-124814484784&-323386252325\\
-377050587335&-448869470372&70859178042&-55846027846&1182451590433\\
585209765500&-541232190289&464444926073&60765711274&-107072278033\\
74541284483&-126095931572&461662797924&455696731011&445244851191\\
48142427836&-53078770252&237226558258&964201483854&-530101168667
\end{pmatrix}
\]
From the 20 matrices $E_{jk}+E_{kj},i(E_{jk}-E_{kj})$, ordered
lexicographically by $j<k$, omit the real directions $(1,3),(2,3),(4,5)$
and the imaginary direction $(2,5)$. Call the remaining matrices
$B_1,\ldots,B_{16}$. On the cube $|x_k|\le10^{-10}$ put
\[
 W=(I-i(A+\sum_k x_kB_k))^{-1},\qquad V(x)=2W-I
\]
Put $F(x)=\Phi(V(x))$ and $H(x)=D^2F(x)$.
This parametrization is exactly unitary. Differentiation uses
$V_j=2iWB_jW$ and
$V_{jk}=-2W(B_jWB_k+B_kWB_j)W$.
The interval computation in \texttt{stalled5.py}, at $40$ decimal digits,
constructs rational matrices $Y,P$ and verifies throughout the cube
\[
 \|Y\nabla F(0)\|_\infty<7\cdot10^{-13},\qquad
 \|I-YH(x)\|_\infty<.077,\qquad
 \|P^tH(x)P-I\|_\infty<.075
\]
Here $Y$ approximates the inverse midpoint Hessian and $P$ the inverse
transpose Cholesky factor, with entries rounded to 16 significant
decimals. The map $x-Y\nabla F(x)$ contracts the cube strictly into
itself, so it has a unique interior fixed point; $H$ is positive definite.
The 16 chart tangents and 9 phase tangents are independent:
for their real $25\times25$ coordinate matrix $T(x)$, a rational
preconditioner $Z$ gives $\|I-ZT(x)\|_\infty<2.2\cdot10^{-5}$.
Thus the chart is transverse to the phase orbits. Finally the same
intervals give $\min|V(x)_{ij}|>.1273$ and the displayed bounds on
$\|R(V(x))\|_F$. Phase invariance proves the asserted strict local
minimum modulo phases.
\end{proof}

\subsection{Certified critical points}\label{criticalpoints}
We finish by validating the critical points behind the catalogue.
The certificate concerns local existence, uniqueness and inertia;
the numerical grouping is not used to infer completeness or pairwise
inequivalence of all $40$ entries.

\begin{thm}\label{criticalcatalogue}
Each of the $40$ catalogue representatives is within $10^{-7}$ in
Frobenius norm of a full-support critical point of the $1$-norm on $U_5$.
This point is locally unique modulo row and column phases, and its
full Hessian has inertia
\[
 (16-p,9,p)
\]
where $p$ is the positive index listed in the table. Its norm rounds
to the displayed value. In particular, the $1$-norm has at least $35$
distinct nonflat critical values. For group $39$, all row and column
pair curvatures are less than $-.22$.
\end{thm}
\begin{proof}
The rational chart data are listed in \texttt{data/critical5-input.json}.
For each representative an integer matrix $C$ defines a self-adjoint
matrix $A$ by $A_{ii}=10^{-20}C_{ii}$ and
$A_{ij}=10^{-20}(C_{ij}+iC_{ji})$ for $i<j$.
The record selects $16$ matrices $B_a$ from the canonical basis of
$5$ diagonal matrices, followed by $E_{jk}+E_{kj},i(E_{jk}-E_{kj})$
for $j<k$ in lexicographic order. Its indices start at $1$.
On the cube $|x_a|\le r=10^{-12}$ put
\[
 W(x)=\left(I-i\left(A+\sum_a x_aB_a\right)\right)^{-1},\qquad
 V(x)=2W(x)-I,\qquad f(x)=L(V(x))
\]
This is an exactly unitary chart. The record also gives an integer
$k\in[-16,16]$ and the unit scalar $z=(4+ik)/(4-ik)$.
Interval bounds throughout the cube give
\[
 \min_{ij}|V(x)_{ij}|>.016,\qquad
 f(x)<11.1<5\sqrt5,\qquad
 \|V(x)/z-U^{(0)}\|_F<10^{-7}
\]
Here $U^{(0)}$ is the corresponding stored representative. Thus all
points in the chart have full support and are nonflat.

The derivatives used by \texttt{critical5.py} are
\[
 V_a=2iWB_aW,\qquad
 V_{ab}=-2W(B_aWB_b+B_bWB_a)W
\]
Writing $m_{ij}=|V_{ij}|$ and
$t_{a,ij}=Re(\overline{V_{ij}}(V_a)_{ij})$, these give
\[
 f_a=\sum_{ij}\frac{t_{a,ij}}{m_{ij}},\qquad
 f_{ab}=\sum_{ij}\left(
 \frac{Re\big(\overline{(V_a)_{ij}}(V_b)_{ij}
                   +\overline{V_{ij}}(V_{ab})_{ij}\big)}{m_{ij}}
 -\frac{t_{a,ij}t_{b,ij}}{m_{ij}^3}\right)
\]
Put $H(x)=D^2f(x)$. At $50$ decimal digits, interval arithmetic
verifies the following bounds for all $40$ charts:
\[
 \|Y\nabla f(0)\|_\infty<10^{-20},\qquad
 \|I-YH(x)\|_\infty<.002,\qquad
 \|P^tH(x)P-D\|_\infty<.001
\]
Here $Y,P$ are rational matrices, proposed by midpoint inversion and
diagonalization and rounded to $17$ significant digits. The diagonal
matrix $D$ has $16-p$ entries $-1$ and $p$ entries $1$.
The first 2 bounds make $x-Y\nabla f(x)$ a contraction of the cube
strictly into itself, with a unique zero of $\nabla f$.
The third fixes the inertia of $H$ by congruence and a nonsingular
homotopy to $D$; numerical eigenvalue signs alone are not used as proof.

To pass to full criticality, form the real $25\times25$ coordinate
matrix $T(x)$ of the $16$ self-adjoint tangent generators $-iV^*V_a$,
the $5$ generators $E_{jj}$, and the $4$ generators $V^*E_{jj}V$
with $j<5$. Coordinates are diagonal entries followed by real and
imaginary parts of upper entries. A rational preconditioner gives
\[
 \|I-ZT(x)\|_\infty<10^{-6}
\]
Thus the chart is transverse to the phase orbits. Phase invariance
turns its critical point into a critical point on $U_5$, locally unique
modulo phases. At that point the $9$ phase directions lie in the full
Hessian kernel, while its restriction to the chart has the certified
nondegenerate inertia. This proves the stated full inertia.

Finally, the norm enclosure for each chart lies strictly within
$5\cdot10^{-7}$ of the displayed value. The table contains $35$
different such values, whose open rounding intervals are disjoint.
For group $39$, interval evaluation of Lemma~\ref{twocol} on $V(x)$
and its transpose bounds all $20$ pair curvatures above by $-.22$.
The individual checks and enclosures are recorded in
\texttt{data/critical5.json} and reproduced by \texttt{critical5.py}.
\end{proof}

\phantomsection\addcontentsline{toc}{section}{References}
\begin{thebibliography}{99}
\setlength{\itemsep}{2pt}
\setlength{\parsep}{0pt}
\bibitem{AMS} P.-A. Absil, R. Mahony and R. Sepulchre,
Optimization Algorithms on Matrix Manifolds, Princeton University
Press (2008).
\bibitem{AMM} G. Auberson, A. Martin and G. Mennessier,
On the reconstruction of a unitary matrix from its moduli,
\emph{Commun. Math. Phys.} \textbf{140} (1991), 417--437.
\bibitem{BCS} T. Banica, B. Collins and J.-M. Schlenker,
On orthogonal matrices maximizing the $1$-norm,
\emph{Indiana Univ. Math. J.} \textbf{59} (2010), 839--856.
\bibitem{BNp} T. Banica and I. Nechita, Almost Hadamard matrices:
the case of arbitrary exponents, \emph{Discrete Appl. Math.}
\textbf{161} (2013), 2367--2379.
\bibitem{BN} T. Banica and I. Nechita, Almost Hadamard matrices with
complex entries, \emph{Adv. Oper. Theory} \textbf{3} (2018), 137--177.
\bibitem{BNS} T. Banica, I. Nechita and J.-M. Schlenker,
Submatrices of Hadamard matrices: complementation results,
\emph{Electron. J. Linear Algebra} \textbf{27} (2014), 197--212.
\bibitem{BNZ} T. Banica, I. Nechita and K. \.Zyczkowski,
Almost Hadamard matrices: general theory and examples,
\emph{Open Syst. Inf. Dyn.} \textbf{19} (2012), 1250024.
\bibitem{BKP} N. Bebiano, A. Kova\v{c}ec and J. da Provid\^encia,
The validity of the Marcus--de Oliveira conjecture for essentially
Hermitian matrices, \emph{Linear Algebra Appl.} \textbf{197--198}
(1994), 411--427.
\bibitem{BMP} N. Bebiano, J.K. Merikoski and J. da Provid\^encia,
On a conjecture of G.N. de Oliveira on determinants,
\emph{Linear Multilinear Algebra} \textbf{20} (1987), 167--170.
\bibitem{BebN} N. Bebiano and H. Nakazato, Critical loci of the $4\times4$
unistochastic map, DMUC Preprint 25-03 (2025).
\bibitem{BP} N. Bebiano and J.P. da Provid\^encia,
Revisiting the Marcus--de Oliveira conjecture,
\emph{Mathematics} \textbf{13} (2025), 711.
\bibitem{BEKTZ} I. Bengtsson, \AA. Ericsson, M. Ku\'s, W. Tadej and
K. \.Zyczkowski, Birkhoff's polytope and unistochastic matrices,
$N=3$ and $N=4$, \emph{Commun. Math. Phys.} \textbf{259} (2005), 307--324.
\bibitem{BW} S. Brierley and S. Weigert, Maximal sets of mutually unbiased
quantum states in dimension 6, \emph{Phys. Rev. A} \textbf{78} (2008), 042312.
\bibitem{D} S.W. Drury, OMC for scalar multiples of unitaries,
\emph{Linear Algebra Appl.} \textbf{422} (2007), 318--325.
\bibitem{DZ} C. Dunkl and K. \.Zyczkowski, Volume of the set of
unistochastic matrices of order 3 and the mean Jarlskog invariant,
\emph{J. Math. Phys.} \textbf{50} (2009), 123521.
\bibitem{F} M. Fiedler, Bounds for the determinant of the sum of
Hermitian matrices, \emph{Proc. Amer. Math. Soc.} \textbf{30} (1971), 27--31.
\bibitem{M} M. Marcus, Derivations, Pl\"ucker relations, and the numerical
range, \emph{Indiana Univ. Math. J.} \textbf{22} (1973), 1137--1149.
\bibitem{MM} J.I. Mulero-Mart\'inez, A variational framework for
determinantal inequalities of normal matrices: successes and obstructions,
\emph{Linear Algebra Appl.} \textbf{740} (2026), 19--38.
\bibitem{O} G.N. de Oliveira, Normal matrices (research problem),
\emph{Linear Multilinear Algebra} \textbf{12} (1982), 153--154.
\bibitem{PB} J. da Provid\^encia and N. Bebiano, Matrices satisfying a
conjecture of G.N. de Oliveira on determinants,
\emph{Linear Algebra Appl.} \textbf{78} (1986), 187--198.
\bibitem{RGZ} G. Rajchel, A. G\k{a}siorowski and K. \.Zyczkowski,
Robust Hadamard matrices, unistochastic rays in Birkhoff polytope and
equi-entangled bases in composite spaces,
\emph{Math. Comput. Sci.} \textbf{12} (2018), 473--490.
\bibitem{RKKZ} G. Rajchel-Mieldzio\'c, K. Korzekwa, Z. Pucha\l a and
K. \.Zyczkowski, Algebraic and geometric structures inside the Birkhoff
polytope, \emph{J. Math. Phys.} \textbf{63} (2022), 012202.
\bibitem{R} K. Rodtes, A class of normal dilation matrices affirming the
Marcus--de Oliveira conjecture, \emph{Oper. Matrices} \textbf{15} (2021), 127--130.
\bibitem{TZguide} W. Tadej and K. \.Zyczkowski,
A concise guide to complex Hadamard matrices,
\emph{Open Syst. Inf. Dyn.} \textbf{13} (2006), 133--177.
\bibitem{TZ} W. Tadej and K. \.Zyczkowski, Defect of a unitary matrix,
\emph{Linear Algebra Appl.} \textbf{429} (2008), 447--481.
\bibitem{T} G. Teschl, Ordinary Differential Equations and
Dynamical Systems, Graduate Studies in Mathematics \textbf{140},
American Mathematical Society (2012).
\bibitem{ZKSS} K. \.Zyczkowski, M. Ku\'s, W. S\l omczy\'nski and
H.-J. Sommers, Random unistochastic matrices,
\emph{J. Phys. A} \textbf{36} (2003), 3425--3450.
\end{thebibliography}
\end{document}
